\documentclass[journal]{IEEEtran}
\usepackage{amsmath,amssymb,mathtools}
\usepackage{amsthm}
\usepackage{booktabs,array}
\usepackage{algorithm}
\usepackage{algpseudocode}
\usepackage{enumitem}
\usepackage{microtype}
\usepackage[hidelinks]{hyperref}
\newtheoremstyle{plainn}{}{}{\itshape}{}{\bfseries}{.}{.5em}%
  {\thmname{#1}\thmnumber{ #2}\thmnote{ (#3)}}
\theoremstyle{plainn}
\newtheorem{thm}{Theorem}
\newtheorem{lem}[thm]{Lemma}
\newtheorem{prop}[thm]{Proposition}
\newtheorem{cor}[thm]{Corollary}
\theoremstyle{remark}
\newtheorem{remark}{Remark}
\renewcommand{\qedsymbol}{QED}

\begin{document}
\title{Parameterized algorithms for functional observer feasibility}
\author{}
\maketitle

\begin{abstract}
We study the parameterized complexity of deciding whether a functional observer of prescribed
order exists for a linear system with rational data and a fixed block of observer rows. The two
feasibility conditions are coupled, because the observer matrix and its product with the system
matrix are built from the same free rows; restricting the row functionals to the common kernel of
the output matrix and its product with the system matrix removes exactly the directions those two
matrices already account for. Exploiting this coupling, we reduce the free rows to a subspace of a
compressed row-pair space. The reduction identifies the parameter that governs the cost, namely
$\ell=\operatorname{rank}(AN)\le k$, where $N$ is a basis of $\ker[C;CA]$: it is an invariant of
the system and of the fixed functionals alone and does not involve the number $r$ of free rows,
whereas the coordinate count $\kappa_{\rm obs}=r\operatorname{rank}[N\ AN]$ that a chart
parameterization suggests does. On this reduction we build a finite algebraic certificate for the
spectral condition and compress it by Cauchy--Binet to a Gram matrix whose size is governed by
$\ell$ rather than by the number of maximal minors, so no enumeration of maximal minors is required;
charts enter only the fixed-parameter algorithm, where their number is bounded by $4^{\ell}$. This yields a deterministic algorithm with bit complexity $\mathcal L^{O(q^2)}$ for
every fixed order $q$, where $\mathcal L$ is the input bit length; a decision algorithm with bound
$\exp\bigl(O(\ell^2\log(\ell+2))\bigr)\operatorname{poly}(\mathcal L)$, that is, fixed-parameter
tractable in a parameter independent of $r$; monotonicity and saturation of the free-row count,
which make the threshold problem a single decision instance and give the minimum order by binary
search; a construction that returns the observer in the arithmetic model, its bit complexity being
left open; and an exact polynomial-time test for the case of no free rows, based on polynomial left
inverses. The decision problem and its finite threshold version lie in the existential theory of
the reals. These are parameterized upper bounds: they do not resolve the unrestricted complexity,
and the doubly exponential dependence on $\ell$ is structural rather than practical.
\end{abstract}

\begin{IEEEkeywords}
functional observer; minimum order; parameterized complexity; real algebraic geometry; polynomial matrix; fixed-parameter tractability.
\end{IEEEkeywords}

\section{Introduction}\label{sec:intro}

Estimating a linear function of the state, rather than the whole state, is often what a
control loop actually needs. Functional observers do exactly this, and their appeal is
quantitative: the order of a functional observer can be far below the order of a
Luenberger observer for the same system, because only the directions that the functional
sees have to be estimated. The design question is therefore an order question: given the
system and the functionals to be estimated, what is the smallest number of observer rows
for which a functional observer exists?

For a linear system with measured output, Fernando, Trinh and Jennings \cite{fer10}
reduced that question to two rank conditions on the observer matrix: with the observer
rows written as a block $L$ that contains the given functional rows $L_0$, one condition
is algebraic and one is spectral. They characterize the feasible rows for a prescribed
order, but they do not address the cost of deciding whether such rows exist.

The characterization itself has since been completed. A unified algebraic framework is
given in \cite{dar25}. The two recent works \cite{fer26a,fer26b} determine, order by
order, which orders are attainable: \cite{fer26a} identifies through functional
observability indices the minimum dimension permitted by the algebraic condition, the
complete family of augmentations attaining it, and necessary and sufficient spectral
conditions for that dimension to be optimal, while \cite{fer26b} characterizes the
interval between that minimum and the endpoint set by functional observability, using a
nullspace representation of the algebraic condition and a matrix-pencil form of the
spectral one. The latter work also describes an order-by-order procedure: at each order
it parameterizes the nullspace over Grassmann charts with $d(m_L-d)=(n-p-q)(q-r)$ free
parameters per chart and tests the two conditions chart by chart.

Neither work states a running-time or a bit-complexity bound, and the order-by-order
procedure is not by itself an algorithm with a complexity guarantee. Two counts show why. The
number of parameters per chart, $d(m_L-d)=(n-p-q)(q-r)$, grows linearly in $n$ even when the
prescribed order $q$ and the number of free rows $r$ are held fixed, and the number of charts
of the relevant Grassmannian grows with $n$ as well. A procedure that enumerates charts and
tests each one therefore does not establish polynomial time for fixed order, and it does not
identify any quantity that could control the cost.

The present paper takes the two conditions as given, proves nothing about them that is already
known, and studies that complexity. Three things are new. First, the conditions are made
finite: the free rows act only through a subspace of a compressed row-pair space, and the
spectral condition is replaced by a finite algebraic certificate of bounded degree, so that no
enumeration of maximal minors is needed, and the fixed-order algorithm needs no chart enumeration
either; the fixed-parameter algorithm does enumerate charts, but at most $4^{\ell}$ of them. Second, a parameter that
controls the cost is identified, namely $\ell=\operatorname{rank}(AN)$, which is an invariant
of the system and of the fixed functionals alone and does not involve $r$; the parameter that
the chart parameterization suggests, $r\,d_0$, does involve $r$ and can be quadratic in $n$
even when $\ell$ is small. Third, the cost is bounded in the bit model: polynomial in the input
length for fixed order, and fixed-parameter in $\ell$. The characterization of
\cite{fer10,dar25,fer26a,fer26b} supplies the starting point throughout; what is added is the
reduction that makes the conditions finite, the certificate that replaces the chart-by-chart
spectral test, and the analysis of the resulting cost.

Adjacent design problems are known to be hard, which makes the question natural but does
not answer it. Choosing a minimal set of sensors that renders a system functionally
observable, or structurally functionally observable, is NP-hard \cite{zha25}, and
constrained sensor-placement variants are NP-complete \cite{dey21}. Static output
feedback with a spectral predicate is NP-hard \cite{lof26,bar26} and, in its real
feasibility form, $\exists\mathbb R$-complete \cite{liu26}; affine rank minimization is
$\exists\mathbb R$-complete already at low rank \cite{maj26}. None of these results
transfers. In sensor placement the quantified object is a subset of the rows of the
output matrix, so the output matrix is itself a design variable, whereas here it is fixed
and the design variable is the set of observer rows. The predicates also differ: a
condition on the spectral abscissa of a closed loop, against a constant-rank condition
that has to hold for every complex value of the spectral variable. A reduction would
moreover have to preserve the single block of free rows that occurs in both conditions,
together with the common origin of the compressed matrices, whereas rank-constrained
formulations leave every entry of the constrained matrix free.

The approach is to make the decision problem finite-dimensional. Restricting row
functionals to the common kernel of the output matrix and its product with the system
matrix removes exactly the directions that those two matrices already account for, and
turns the two conditions into a row-space inclusion and into the constant rank of a
pencil. A free observer row is then determined by its coordinates in the space of
attainable row pairs, and the $r$ free rows are seen to act only through the subspace $W$
they span there. That reading gives two parameters: the dimension $\kappa_{\rm obs}=rd_0$
of the coordinates of the free rows (Section~\ref{sec:dimension}), and the dimension
$\ell=\operatorname{rank}(AN)\le k$ of the space in which the feasible $W$ live, which does
not involve $r$ (Section~\ref{sec:subspace}). Three consequences are developed. The spectral
condition becomes a finite algebraic certificate compressed to a size governed by $\ell$
(Section~\ref{sec:certificate}); deciding feasibility becomes a real-algebraic problem
whose number of variables is controlled by either the prescribed order or by $\ell$
(Sections~\ref{sec:fixedq} and~\ref{sec:fpt}); and the saturation of the free-row count at
$d-s_0$ turns the threshold problem into a single decision instance and the minimum-order
problem into a binary search.

\subsection{Contributions and organization}\label{sec:contrib}

The contributions are the following.

\begin{enumerate}
\item A compression of the two observer conditions, which are taken from the existing
characterization, to $O(q^2)$ variables independently of the state dimension, and a
reduction of the free rows to a subspace $W$ of the compressed row-pair space, in which the
two conditions are a row-space inclusion and the constant rank of a pencil.
\item The identification of the parameter $\ell=\operatorname{rank}(AN)\le k$ that governs
the complexity bound. It is an invariant of the system and the fixed rows alone: it does not
involve the number $r$ of free rows, whereas the chart parameterization of \cite{fer26b}
carries $(n-p-q)(q-r)$ parameters per chart for fixed order $q$.
\item A lower bound on the minimum order that comes from the spectral condition rather than
from the algebraic one: $\operatorname{rank}[U_0;V_0]\le\omega\le s_0+t-1$ with $U_0=\operatorname{row}(L_0N)$
and $V_0=\operatorname{row}(L_0AN)$, hence
$q_{\min}\ge r_0+\max(0,\operatorname{rank}[U_0;V_0]-s_0+1)$. It is computed in polynomial time,
it is independent of the algebraic bound $q_0$ of \cite{fer26a,fer26b}, and on the instance of
Appendix~\ref{app:example} it is strictly stronger.
\item A finite algebraic certificate for the spectral condition, and its compression by
Cauchy--Binet to a matrix of size controlled by $\ell$; no enumeration of maximal minors is
required.
\item A decision algorithm for every fixed order, with a bit-complexity bound
$\mathcal L^{O(q^2)}$ in the input length $\mathcal L$ and the prescribed order $q$.
\item A decision algorithm that is fixed-parameter tractable in $\ell$, with bound
$\exp\bigl(O(\ell^2\log(\ell+2))\bigr)\operatorname{poly}(\mathcal L)$, together with the
monotonicity and saturation of the free-row count, which make the threshold problem a single
decision instance and give minimum order by binary search, and an exact polynomial-time test
for the case of no free rows.
\item Membership of the prescribed-order decision problem, and of its finite threshold
version, in the existential theory of the reals.
\end{enumerate}

The paper is organized as follows. Section~\ref{sec:prelim} fixes the notation and states
the decision, optimization and construction problems. Section~\ref{sec:compressed} proves
the compression, Section~\ref{sec:dimension} identifies the effective dimension, and
Section~\ref{sec:certificate} develops the certificate. Sections~\ref{sec:fixedq}
and~\ref{sec:fpt} give the two algorithms, Section~\ref{sec:existsr} records the
existential-real membership, and Section~\ref{sec:conclusions} concludes.

\section{Preliminaries}\label{sec:prelim}

\subsection{Notation and conventions}\label{sec:notation}

Consider the linear time-invariant system
\begin{equation}\label{system}
\dot{x}=Ax+Bu,\quad y=Cx,\quad z=L_0x,
\end{equation}
where $A\in\mathbb Q^{n\times n}$, $C\in\mathbb Q^{p\times n}$ and
$L_0\in\mathbb Q^{r_0\times n}$.

Let $q\ge r_0$ be the prescribed observer order, with $q\le n$, and set $r=q-r_0$. The
observer rows are the rows of the block
\begin{equation*}
L=\begin{bmatrix}L_0\\R\end{bmatrix},\qquad R\in\mathbb R^{r\times n},
\end{equation*}
where $L_0$ is fixed while $R$ collects the $r$ free rows; only real observers are
considered.

All given rational data are encoded in binary, and $\mathcal L$ denotes the total input
bit length. The free parameter $R$ is not part of the given input data and is therefore
not counted in $\mathcal L$. State vectors are column vectors, and matrices act on them by
the usual left multiplication.

Let $P(\lambda)$ be a rational matrix pencil of generic rank $\omega$, and call its
$\omega\times\omega$ minors maximal minors. Ranks at spectral points $\lambda\in\mathbb C$
are always computed over $\mathbb C$. Because the maximal minors have rational
coefficients, their gcd is a nonzero constant in $\mathbb C[\lambda]$ exactly when it is a
nonzero constant in $\mathbb Q[\lambda]$; this equivalence is used without further
mention.

\subsection{The observer conditions}\label{sec:conditions}

A functional observer of order $q$ is a matrix $L=\begin{bmatrix}L_0 \\ R\end{bmatrix}\in\mathbb R^{q\times n}$ such
that
\begin{equation}\label{eq:5a}
\operatorname{rank}\begin{bmatrix}LA\\ L\\ CA\\ C\end{bmatrix}
=\operatorname{rank}\begin{bmatrix}L\\ CA\\ C\end{bmatrix},
\end{equation}
\begin{equation}\label{eq:5b}
\operatorname{rank}\begin{bmatrix}\lambda L-LA\\ CA\\ C\end{bmatrix}
=\operatorname{rank}\begin{bmatrix}L\\ CA\\ C\end{bmatrix},
\quad \forall\lambda\in\mathbb C.
\end{equation}
These are the Fernando--Trinh--Jennings conditions, numbered (5a) and (5b) in
\cite{fer10}. Condition~\eqref{eq:5a} is algebraic; \eqref{eq:5b} is spectral, and it is
required at every $\lambda\in\mathbb C$, in particular at all finite eigenvalues of $A$.

\subsection{Decision, optimization and construction}\label{sec:problems}

Three problems are posed for \eqref{eq:5a} and \eqref{eq:5b}, in increasing order of
demand.

\emph{Decision.} $\mathrm{DEC}$ takes $(A,C,L_0)$ and $q$ and asks whether some $R$
satisfies \eqref{eq:5a} and \eqref{eq:5b}; no witness is required. This is the problem
studied in this paper.

\emph{Optimization.} $\mathrm{OPT}$ asks for the minimum feasible order
\begin{equation*}
q_{\min}=\min\{q:\mathrm{DEC}(q)\},
\end{equation*}
where $\mathrm{DEC}(q)$ abbreviates $\mathrm{DEC}(A,C,L_0,q)$.

\emph{Construction.} $\mathrm{FIND}$ asks for a witness $R$ at the minimum feasible
order, $q=q_{\min}$.

The threshold problem is
\begin{equation}\label{eq:mindec}
\mathrm{MINDEC}(Q)=\bigvee_{q=r_0}^{\min(Q,n)}\mathrm{DEC}(q),
\end{equation}
a finite disjunction of decision instances. The algorithms below decide $\mathrm{DEC}$;
their affirmative certificates contain a witness. Section~\ref{sec:dimension} also shows
that $\mathrm{DEC}$ is monotone in $q$ and saturates at $r_0+m$, which reduces
\eqref{eq:mindec} to a single instance and turns $\mathrm{OPT}$ into a search over at most
$2\ell+1$ values; $\mathrm{FIND}$ is treated in Section~\ref{sec:fpt}, where the witness is
returned as exact real algebraic numbers.

\begin{table}[!t]
\centering
\caption{Notation. The first block is fixed by the input, the second by the compression, the third by
the subspace reduction of Section~\ref{sec:subspace}.}
\label{tab:notation}
\footnotesize
\begin{tabular}{@{}ll@{}}
\toprule
Symbol & Meaning \\
\midrule
$K_0=\ker[C;CA]$ & common kernel of $C$ and $CA$ \\
$k=\dim K_0$ & its dimension \\
$N$ & basis matrix of $K_0$ \\
$M=[N\ AN]$ & compressed row-pair matrix \\
\midrule
$H=LN$, $H_A=LAN$ & compressed observer blocks \\
$\omega=\operatorname{rank}H$ & compressed order \\
$P(\lambda)=\lambda H-H_A$ & compressed pencil \\
\midrule
$\mathcal T=\operatorname{row}M$ & attainable row-pair space \\
$d_0=\dim\mathcal T$ & its dimension \\
$\kappa_{\rm obs}=r\,d_0$ & number of free-row coordinates \\
$S=\operatorname{row}(AN)\subseteq\mathbb R^k$ & space of the second blocks \\
$\ell=\operatorname{rank}(AN)\le k$ & its dimension; the FPT parameter \\
$B=LM=[H\ H_A]$ & compressed row pair \\
$W=\operatorname{row}B$ & subspace spanned by the row pairs \\
$W_0=\operatorname{row}(L_0M)$ & contribution of the fixed rows \\
$s_0=\dim W_0$, $m=d_0-s_0$ & \\
$\mathcal V=\mathcal T\cap(S\times S)$ & admissible ambient space \\
$m'=\dim\mathcal V-s_0$ & chart dimension \\
$\pi_1,\pi_2,\pi_\lambda$ & projections of $\mathbb R^{2k}$ \\
\bottomrule
\end{tabular}
\end{table}

Two published bounds confine the orders that have to be examined. Let $q_0$ be the algebraic
lower bound of \cite{fer26a,fer26b}: the smallest order at which \eqref{eq:5a} can be satisfied,
computed from the functional observability indices by a greedy rank recursion, hence in polynomial
time. Every feasible order satisfies
\begin{equation}\label{eq:range}
q_0\le q\le n_0-p,\qquad n_0=\operatorname{rank}\bigl[C;CA;\dots;CA^{n-1}\bigr],
\end{equation}
see \cite{fer26a,fer26b}; for an observable pair $n_0=n$. The upper end is not merely a bound in the
row-slot model: by Lemma~\ref{lem:monotone} the predicate is monotone in $q$, so every order at or
above $n_0-p$ is feasible as soon as that one is, and no information is lost by stopping there. In
terms of the quotient dimension the range \eqref{eq:range} reads
$q_0-r_0\le t\le n_0-p-r_0$, and the reduction of Section~\ref{sec:subspace} adds the independent
bound $t\le2\ell$ from Proposition~\ref{prop:minorder}.

\section{The compressed pencil}\label{sec:compressed}

Both observer conditions depend on the free rows only through the products $LN$ and $LAN$, where $N$ is a basis matrix of the common kernel of $C$ and $CA$. This section performs that compression and rewrites \eqref{eq:5a} and \eqref{eq:5b} in terms of the compressed data.

\subsection{The common kernel and the compression}\label{sec:compress}

Let
\begin{equation*}
K_0=\ker\begin{bmatrix}C\\ CA\end{bmatrix},\qquad k=\dim K_0,
\end{equation*}
and let $N\in\mathbb R^{n\times k}$ be a basis matrix for $K_0$. Define
\begin{equation*}
H=LN\in\mathbb R^{q\times k},\quad H_A=LAN\in\mathbb R^{q\times k},\quad \omega=\operatorname{rank}H .
\end{equation*}
The compressions below rest on a single rank identity.

\begin{lem}[Rank identity]\label{lem:rankid}
Let $P\in\mathbb C^{m\times n}$ and let $N_P\in\mathbb C^{n\times r}$ be a basis matrix for $\ker P$. Then, for every $M\in\mathbb C^{p\times n}$,
\begin{equation}\label{eq:rankid}
\operatorname{rank}\begin{bmatrix}P\\ M\end{bmatrix}
=\operatorname{rank}P+\operatorname{rank}(MN_P).
\end{equation}
\end{lem}

\begin{proof}
Let $f(x)=xN_P$. Since $\operatorname{col}N_P=\ker P$, the kernel of $f$ is $\operatorname{row}P$,
and $\operatorname{row}(MN_P)$ is the image of the restriction of $f$ to $\operatorname{row}M$, so
rank--nullity gives $\operatorname{rank}[P;M]=\operatorname{rank}P+\operatorname{rank}(MN_P)$; this
is the standard rank identity for a stacked matrix \cite[Ch.~0]{hj13}.
\end{proof}

\subsection{The compression and the pencil criterion}\label{sec:compression}

\begin{lem}[Compression and pencil criterion]\label{lem:compression}
With $H$, $H_A$ and $\omega$ as above, the following statements hold.

(i) Condition \eqref{eq:5a} holds if and only if $\operatorname{row}H_A\subseteq\operatorname{row}H$.

(ii) If \eqref{eq:5a} holds, then there exists $F\in\mathbb R^{q\times q}$ such that $H_A=FH$.

(iii) Assume \eqref{eq:5a}. Then \eqref{eq:5b} holds if and only if
\begin{equation*}
\operatorname{rank}(\lambda H-H_A)=\omega,\qquad \forall\lambda\in\mathbb C .
\end{equation*}
Equivalently, the $\omega\times\omega$ maximal minors of $P(\lambda)=\lambda H-H_A$ have no common complex root; equivalently, their greatest common divisor in $\mathbb C[\lambda]$ is a nonzero constant.

(iv) If \eqref{eq:5a} and \eqref{eq:5b} hold, then $\operatorname{rank}H_A=\operatorname{rank}H$.
\end{lem}

\begin{proof}
Set $P_0=\begin{bmatrix}C\\CA\end{bmatrix}$. Since $K_0=\ker P_0$, the matrix $N$ is a basis matrix for $\ker P_0$.

\emph{(i)} The rows of $P_0$ occur among the rows of both matrices below, so Lemma~\ref{lem:rankid} with $P=P_0$ and with $M=\begin{bmatrix}LA\\L\end{bmatrix}$, $M=L$ gives
\begin{equation*}
\begin{aligned}
\operatorname{rank}\begin{bmatrix}LA\\ L\\ CA\\ C\end{bmatrix}
&=\operatorname{rank}P_0+\operatorname{rank}\begin{bmatrix}H_A\\ H\end{bmatrix},\\
\operatorname{rank}\begin{bmatrix}L\\ CA\\ C\end{bmatrix}
&=\operatorname{rank}P_0+\operatorname{rank}H .
\end{aligned}
\end{equation*}
Up to row permutations the left-hand sides are the two sides of \eqref{eq:5a}, so
\begin{equation*}
\eqref{eq:5a}
\iff
\operatorname{rank}\begin{bmatrix}H_A\\ H\end{bmatrix}=\operatorname{rank}H ,
\end{equation*}
which is equivalent to
\begin{equation}\label{eq:5a-1}
\operatorname{row}H_A\subseteq\operatorname{row}H .
\end{equation}

\emph{(ii)} By \eqref{eq:5a-1} every row of $H_A$ is a linear combination of the rows of $H$. Let $F$ have as its $i$-th row the coefficient vector of the $i$-th row of $H_A$; then
\begin{equation}\label{eq:5a-2}
H_A=FH .
\end{equation}

\emph{(iii)} Set
\begin{equation*}
M(\lambda)=\begin{bmatrix}
	\lambda L-LA\\CA\\C
\end{bmatrix},\qquad J=\begin{bmatrix}
L \\ CA \\C
\end{bmatrix}.
\end{equation*}
Since $N$ is a basis matrix for $K_0$,
\begin{equation*}
M(\lambda)N=\begin{bmatrix}
	\lambda H-H_A \\ 0 \\ 0
\end{bmatrix}=\begin{bmatrix}
P(\lambda) \\ 0 \\ 0
\end{bmatrix},\qquad JN=\begin{bmatrix}
H \\ 0 \\ 0
\end{bmatrix} .
\end{equation*}
Since the row space of $P_0$ is a subspace of the row spaces of both $M(\lambda)$ and $J$,
\begin{equation*}
\operatorname{rank}\begin{bmatrix}P_0\\ M(\lambda)\end{bmatrix}=\operatorname{rank}M(\lambda),
\qquad
\operatorname{rank}\begin{bmatrix}P_0\\ J\end{bmatrix}=\operatorname{rank}J .
\end{equation*}
So Lemma~\ref{lem:rankid} with $P=P_0$ yields
\begin{equation*}
\begin{aligned}
\operatorname{rank}M(\lambda)&=\operatorname{rank}P_0+\operatorname{rank}P(\lambda),\\
\operatorname{rank}J&=\operatorname{rank}P_0+\operatorname{rank}H .
\end{aligned}
\end{equation*}
Hence \eqref{eq:5b} is equivalent to
\begin{equation}\label{eq:5b-1}
\operatorname{rank}P(\lambda)=\operatorname{rank}H=\omega,\ \forall \lambda\in\mathbb C.
\end{equation}
 Assume now \eqref{eq:5a}. By \eqref{eq:5a-2},
\begin{equation*}
P(\lambda)=\lambda H-H_A=(\lambda I-F)H ,
\end{equation*}
so every row of $P(\lambda)$ is a linear combination of the rows of $H$, that is, $\operatorname{row}P(\lambda)\subseteq\operatorname{row}H$, and therefore
\begin{equation}\label{eq:5b-2}
\operatorname{rank}P(\lambda)\le\omega,\qquad \forall\lambda\in\mathbb C .
\end{equation}
Combining \eqref{eq:5b-1} with the bound \eqref{eq:5b-2}, condition \eqref{eq:5b} becomes the requirement that some $\omega\times\omega$ minor of $P(\lambda)$ be nonzero at every $\lambda$:
\begin{equation*}
\eqref{eq:5b}
\iff
\forall\lambda\in\mathbb C\ \ \exists\,\omega\times\omega\text{ minor } f_i:\ f_i(\lambda)\neq0 .
\end{equation*}
A common complex root of the maximal minors is therefore exactly a value of $\lambda$ at which all of them vanish, which proves the two equivalent forms in the statement; for $\omega=0$ the condition is vacuous.

\emph{(iv)} Assume \eqref{eq:5a} and \eqref{eq:5b}. By (iii), \eqref{eq:5b-1} holds for every $\lambda$; at $\lambda=0$, where $P(0)=-H_A$,
\begin{equation*}
\operatorname{rank}H_A=\operatorname{rank}P(0)=\operatorname{rank}H . \qedhere
\end{equation*}
\end{proof}

\section{The effective free-row dimension}\label{sec:dimension}

\subsection{The space of attainable row pairs}\label{sec:free}

Each free row $r_i\in\mathbb R^{1\times n}$ of $R$ determines the row pair
\begin{equation*}
(H_i,(H_A)_i)=r_i\,[N\ \ AN] \in\mathbb R^{1\times 2k}.
\end{equation*}
Consequently, this attainable pair is uniquely specified by its coordinates relative to a basis of the subspace of attainable pairs
\begin{equation*}
\begin{gathered}
\mathcal T=\operatorname{row}[N\ \ AN]\subseteq\mathbb R^{1\times 2k},\\
d_0=\dim\mathcal T=\operatorname{rank}[N\ \ AN].
\end{gathered}
\end{equation*}
Therefore, a rational basis of $\mathcal T$ allows each attainable pair generated by a free row to be parameterized by exactly $d_0$ real coordinates.

\begin{lem}[Attainable row pairs and the effective dimension]\label{lem:attainable}
Every free row pair lies in $\mathcal T$, and writing the $r$ free rows in a basis of $\mathcal T$ represents them by
\begin{equation}\label{eq:kappa}
\kappa_{\rm obs}=r d_0
\end{equation}
coordinates. Moreover, when $r\ge1$,
\begin{equation}\label{lem3-2}
k\le d_0\le rd_0=\kappa_{\rm obs}.
\end{equation}
\end{lem}

\begin{proof}
The identity $(H_i,(H_A)_i)=r_i[N\ AN]$ places every free row pair in $\mathcal T$. Consider the linear map
\begin{equation*}
\Phi:\mathbb R^{1\times n}\to\mathcal T,\qquad \Phi(x)=x\,[N\ \ AN].
\end{equation*}
Its matrix relative to the standard basis of $\mathbb R^{1\times n}$ and a basis of $\mathcal T$ is $[N\ AN]$, so $\operatorname{rank}\Phi=\operatorname{rank}[N\ AN]$ and rank--nullity gives
\begin{equation*}
d_0=\dim\mathcal T=\dim\operatorname{im}\Phi=n-\dim\ker\Phi=\operatorname{rank}[N\ AN].
\end{equation*}
For the $r$ free rows set
\begin{equation*}
\Psi:(\mathbb R^{1\times n})^r\to\mathcal T^{\,r},\quad
\Psi(r_1,\dots,r_r)=\bigl(\Phi(r_1),\dots,\Phi(r_r)\bigr).
\end{equation*}
Its image is $\mathcal T^{\,r}$, so the number of coordinates after parametrization is
\begin{equation*}
\dim\mathcal T^{\,r}=r\dim\mathcal T=rd_0=\kappa_{\rm obs}.
\end{equation*}
Since $N$ has full column rank,
\begin{equation*}
\operatorname{rank}N=k,\qquad \operatorname{rank}[N\ AN]\ge\operatorname{rank}N=k,
\end{equation*}
that is,
\begin{equation}\label{lem3-1}
k\le d_0 .
\end{equation}
Since $r\ge1$, we have $d_0\le rd_0=\kappa_{\rm obs}$, which together with \eqref{lem3-1} is \eqref{lem3-2}.
\end{proof}

\begin{remark}[Basis invariance]
If the basis matrix is changed via $N\mapsto NQ$ for some nonsingular $Q\in\mathbb R^{k\times k}$, then the concatenated matrix transforms as
\begin{equation*}
[N\ \ AN]\mapsto[N\ \ AN]\operatorname{diag}(Q,Q),
\end{equation*}
Since \(\operatorname{diag}(Q,Q)\) is nonsingular, the row space \(\mathcal T\) remains unchanged. Furthermore, to preserve the generated attainable pairs \(r_i[N\ \ AN]\), the free row coordinates must transform via the nonsingular linear map \(r_i \mapsto r_i \operatorname{diag}(Q^{-1}, Q^{-1})\). Consequently, both the dimension \(d_0 = \dim\mathcal T\) and the total number of parameters \(\kappa_{\rm obs} = r d_0\) are invariant with respect to the choice of basis.
\end{remark}

\subsection{The subspace reduction}\label{sec:subspace}

The free rows enter the two conditions only through the row space of the compressed
pair. Writing the row pair of a free row as a row of $M=[N\ AN]$, let
\begin{equation}\label{eq:W-def}
B=LM=[H\ H_A]\in\mathbb R^{q\times2k},\qquad W=\operatorname{row}B\subseteq\mathbb R^{2k},
\end{equation}
so that $W_0=\operatorname{row}(L_0M)\subseteq W\subseteq\mathcal T=\operatorname{row}M$.
Let
\begin{equation}\label{eq:ell}
\begin{gathered}
d=\dim\mathcal T,\qquad s_0=\dim W_0,\qquad m=d-s_0,\\
\ell=\operatorname{rank}(AN)\le k,\qquad S=\operatorname{row}(AN)\subseteq\mathbb R^k ,
\end{gathered}
\end{equation}
and let $\pi_1(h,g)=h$, $\pi_2(h,g)=g$, $\pi_\lambda(h,g)=\lambda h-g$ be the projections of
$\mathbb R^{2k}=\mathbb R^k\times\mathbb R^k$. Because \eqref{eq:5b} quantifies over complex
$\lambda$, the projection $\pi_\lambda$ is applied to the complexification
$W_{\mathbb C}=W\otimes_{\mathbb R}\mathbb C\subseteq\mathbb C^{2k}$, and dimensions are over
$\mathbb C$; a real matrix has the same rank over $\mathbb R$ and over $\mathbb C$.

\begin{lem}[Sandwich and realizability]\label{lem:sandwich}
For $L=[L_0;R]$ one has $W_0\subseteq W\subseteq\mathcal T$ and
$\dim(W/W_0)\le r$. Conversely, every subspace $W$ with these three properties is realized by a
free block $R$ with exactly $t=\dim(W/W_0)$ rows.
\end{lem}

\begin{proof}
$W=\operatorname{row}(L_0M)+\operatorname{row}(RM)$ gives $W\supseteq W_0$ and
$W\subseteq\mathcal T$, and $\dim(W/W_0)\le\operatorname{rank}(RM)\le r$. For the converse, the
map $u\mapsto uM$ is onto $\mathcal T$; choose preimages $u_1,\dots,u_t$ of a basis of
$W/W_0$ and let $R$ consist of those $t$ rows followed by $q-r_0-t$ zero rows. Then
$W_0+\operatorname{row}(RM)=W$.
\end{proof}

\begin{lem}[The conditions depend only on $W$]\label{lem:W-conditions}
Condition \eqref{eq:5a} holds if and only if $\pi_2(W)\subseteq\pi_1(W)$, and condition
\eqref{eq:5b} holds if and only if $\dim\pi_\lambda(W)=\dim\pi_1(W)$ for every
$\lambda\in\mathbb C$.
\end{lem}

\begin{proof}
Both are instances of Lemma~\ref{lem:rankid} with $P=[CA;C]$, whose kernel is $K_0$: for
\eqref{eq:5a} take $Q=[LA;L]$, for \eqref{eq:5b} take $Q=\lambda L-LA$. In the first case
$\operatorname{rank}[LA;L;CA;C]=\operatorname{rank}P+\operatorname{rank}[H_A;H]$ and
$\operatorname{rank}[L;CA;C]=\operatorname{rank}P+\operatorname{rank}H$, so the two are equal
exactly when $\operatorname{row}H_A\subseteq\operatorname{row}H$, which is
$\pi_2(W)\subseteq\pi_1(W)$; in the second, equality reads
$\operatorname{rank}(\lambda H-H_A)=\operatorname{rank}H$.
\end{proof}

\begin{lem}[Structural bound]\label{lem:struct}
Let $W$ be feasible and $\omega=\dim\pi_1(W)>0$. Then
\begin{equation}\label{eq:struct}
\omega+1\le\dim W\le2\omega .
\end{equation}
\end{lem}

\begin{proof}
For the upper bound, $\operatorname{row}(L_0M)\subseteq W$ gives
$W_0\subseteq\mathcal T$, and Lemma~\ref{lem:W-conditions} with $\lambda=0$ gives
$\dim\pi_0(W)=\dim\pi_1(W)=\omega$; since $\pi_0(W)=\pi_2(W)$ up to sign,
$\pi_2(W)=\pi_1(W)=:U$ and $W\subseteq U\times U$, so $\dim W\le2\dim U=2\omega$. For the
lower bound, if $\dim W=\omega$ then $\pi_1|_W$ and $\pi_2|_W$ are isomorphisms and $W$ is the
graph of a linear $F:U\to U$, that is $g=Fh$ on $W$. Over $\mathbb C$, $F$ has an eigenvalue
$\lambda_0$, and $\pi_{\lambda_0}(W)=(\lambda_0I-F)U$ has dimension $<\omega$, contradicting
Lemma~\ref{lem:W-conditions}.
\end{proof}

\begin{lem}[Monotonicity and saturation]\label{lem:monotone}
$\mathrm{DEC}(q)\Rightarrow\mathrm{DEC}(q+1)$, and $\mathrm{DEC}(r_0+r)=\mathrm{DEC}(r_0+m)$
whenever $r\ge m$. Consequently $\mathrm{MINDEC}(Q)=\mathrm{DEC}(\min(Q,n))$ for
$Q\ge r_0$.
\end{lem}

\begin{proof}
Adjoining a zero row to $R$ leaves $H$, $H_A$, and the ranks in \eqref{eq:5a} and
\eqref{eq:5b} unchanged, which gives monotonicity. Saturation follows from
$\dim(W/W_0)\le\dim(\mathcal T/W_0)=m$ and Lemma~\ref{lem:sandwich}: any feasible $W$ with
$\dim(W/W_0)\le r$ is already realized with at most $m$ free rows. The disjunction in
\eqref{eq:mindec} then reduces to its largest term.
\end{proof}

\begin{prop}[Minimum order]\label{prop:minorder}
Let $t_{\min}=\min\{\dim(W/W_0):W \text{ feasible}\}$. Then
$q_{\min}=r_0+t_{\min}$ and $t_{\min}\le2\ell$.
\end{prop}

\begin{proof}
Lemma~\ref{lem:sandwich} identifies feasibility with the existence of a feasible $W$ of
quotient dimension $\le r$, and Lemma~\ref{lem:monotone} shows that the smallest such $r$ is
$t_{\min}$. For the bound, a feasible $W$ satisfies $W\subseteq U\times U$ with
$U\subseteq S$, so $\dim U=\omega\le\ell$ and $\dim W\le2\omega\le2\ell$; hence
$t_{\min}\le\dim W-s_0\le2\ell$.
\end{proof}

\begin{prop}[A computable lower bound]\label{prop:lowerbound}
Let $U_0=\operatorname{row}(L_0N)$ and $V_0=\operatorname{row}(L_0AN)$, and let $W$ be feasible with
$t=\dim(W/W_0)$ and $\omega=\dim\pi_1(W)$. Then
\begin{equation}\label{eq:lb}
\operatorname{rank}\,[\,U_0;\,V_0\,]\ \le\ \omega\ \le\ s_0+t-1,
\end{equation}
hence $t\ge\operatorname{rank}[U_0;V_0]-s_0+1$ and
\begin{equation}\label{eq:qmin-lb}
q_{\min}\ \ge\ r_0+\max\bigl(0,\ \operatorname{rank}[U_0;V_0]-s_0+1\bigr).
\end{equation}
\end{prop}

Proposition~\ref{prop:lowerbound} is the spectral counterpart of the algebraic bound of
\cite{fer26a,fer26b}. That bound, $q_0=\sum_j\eta_j$, limits the orders at which \eqref{eq:5a} can
hold; \eqref{eq:lb} limits the orders at which \eqref{eq:5b} can hold, since it uses only the
condition at $\lambda=0$ and the structural bound. The two are independent, both are computed in
polynomial time, and the larger one applies; on the instance of Appendix~\ref{app:example} the bound
below gives $4$ while $q_0=3$. To our knowledge the literature states no lower bound derived from the
spectral condition: the design conditions of \cite{fer10,dar25,fer26a,fer26b} and the order bound
$q_0$ all concern \eqref{eq:5a}, and the spectral condition is treated there as a test applied at each
order rather than as a source of a bound.

\begin{proof}
Since $W\supseteq W_0$ the projections satisfy $\pi_1(W)\supseteq\pi_1(W_0)=U_0$ and
$\pi_2(W)\supseteq\pi_2(W_0)=V_0$. By Lemma~\ref{lem:W-conditions} at $\lambda=0$ we have
$\pi_1(W)=\pi_2(W)=U$ with $\dim U=\omega$, so $U\supseteq U_0+V_0$ and
$\omega\ge\operatorname{rank}[U_0;V_0]$. The upper bound is Lemma~\ref{lem:struct}:
$\omega+1\le\dim W=s_0+t$. Combining and using $q_{\min}=r_0+t_{\min}$ with $t_{\min}\ge0$ gives
\eqref{eq:qmin-lb}.
\end{proof}

\begin{remark}[On tightness]\label{rem:lbtight}
The bound \eqref{eq:qmin-lb} is attained on the instance of Appendix~\ref{app:example}, where it gives
$4$ and $q_{\min}=4$. We do not claim that it is tight in general: like the algebraic bound $q_0$ it is
a lower bound, and the difference
$q_{\min}-\bigl(r_0+\max(0,\operatorname{rank}[U_0;V_0]-s_0+1)\bigr)$ can be positive. What it does
provide is a test for the orders below it, which then need not be examined at all.
\end{remark}

\begin{remark}[Coordinate change to $S$]\label{rem:YS}
Fix a rational basis of $S$ and extend it to a rational basis of $\mathbb R^k$; let $Y_S$ be
the change-of-basis matrix, so that $X\mapsto XY_S$ puts the coordinates of $S$ first. Right
multiplication by $Y_S$ preserves ranks and the inclusions and equalities between row spaces,
though not row spaces themselves, so Lemma~\ref{lem:W-conditions} and
Lemma~\ref{lem:struct} apply verbatim in the new coordinates and every object involved lives
in $\mathbb R^\ell$. The basis is obtained by column elimination on $AN$, so the entries of
$Y_S$ have bit length $\operatorname{poly}(\mathcal L)$.
\end{remark}

\section{Rank factorization, consistency, and the finite certificate}\label{sec:certificate}

This section derives the finite certificate for the spectral condition
\eqref{eq:5b} in three steps: a pivot factorization of the rank-$\omega$
block, a rank-consistency test for the linear system that the factorization
leaves, and a compression of that test to a Gram matrix whose size is
governed by the parameter $\ell$ of \eqref{eq:ell}.

\subsection{Rank factorization and consistency}\label{sec:consistency}

\paragraph*{Pivot factorization}\label{sec:pivot}
Let $\omega>0$. For a rank-$\omega$ matrix $H$ choose a pivot index set
$J_{p}\subseteq[k]$, $|J_{p}|=\omega$, such that
\begin{equation*}
X=H_{:,J_{p}}\in\mathbb R^{q\times\omega}
\end{equation*}
has full column rank. Then
\begin{equation*}
H=XB,\qquad B\in\mathbb R^{\omega\times k},\qquad B_{:,J_p}=I_\omega .
\end{equation*}
If $H_A=FH$, put $G=FX\in\mathbb R^{q\times\omega}$.
Then
\begin{equation*}
H_A=GB,\qquad P(\lambda)=(\lambda X-G)B .
\end{equation*}
Conversely, if $H_A=GB$ and $X_LX=I_\omega$, where $X_L$ denotes a left
inverse of $X$, then $F=GX_L$ satisfies
\begin{equation*}
FH=GX_LXB=GB=H_A .
\end{equation*}

\paragraph*{Parametric linear system}\label{sec:param}
Let $D$ be a rational matrix with $\ker D=\mathcal T$. For any fixed row index 
$i$, the parameter matrix $B$ is subject to the target constraints
\begin{equation*}
X_iB=h_i^0\in\mathbb R^{1\times k},\qquad G_iB=g_i^0\in\mathbb R^{1\times k},
\end{equation*}
as well as the structural constraint
\begin{equation*}
D\begin{bmatrix}(X_iB)^{\mathsf T}\\ (G_iB)^{\mathsf T}\end{bmatrix}=0.
\end{equation*}
which ensures that the generated row pairs belong to the attainable subspace $\mathcal T$.

Given the pivot factorization condition $B_{:,J_p}=I_\omega$, the known entries of 
$B$ can be eliminated. Let $b$ denote the vectorization of the remaining unknown entries of $B$. The above constraints then reduce to a standard linear system:


\begin{equation*}
M(X,G)b=c(X,G).
\end{equation*}

\begin{lem}[Rank consistency]\label{lem:rankcons}
	For real matrices $M,c$,
	\[
	\exists b:\ Mb=c
	\]
	if and only if
	\[
	\operatorname{rank}(MM^{\mathsf T})
	=\operatorname{rank}(MM^{\mathsf T}+cc^{\mathsf T}).
	\]
	Equivalently, with $\Gamma=MM^{\mathsf T}$ and
	$\Gamma'=MM^{\mathsf T}+cc^{\mathsf T}$,
	\[
	\bigvee_{t=0}^{d_1}
	\left[e_t(\Gamma)>0\wedge e_{t+1}(\Gamma)=0\wedge e_{t+1}(\Gamma')=0\right],
	\]
	where $d_1=\min(\text{rows},\text{cols})$ of $M$ and $e_j$ is the
	$j$-th elementary symmetric polynomial of the eigenvalues.
\end{lem}

\begin{proof}
The equation $Mb=c$ is consistent if and only if $c\in\operatorname{im}M$, which is the classical
consistency criterion for linear systems, $\operatorname{rank}[M\ c]=\operatorname{rank}M$
\cite[Ch.~0]{hj13}. Since
\[
\begin{gathered}
\operatorname{im}[M\ c]=\operatorname{im}M+\operatorname{span}\{c\},\qquad
\operatorname{im}(MM^{\mathsf T})=\operatorname{im}M,\\
\operatorname{im}(MM^{\mathsf T}+cc^{\mathsf T})=\operatorname{im}[M\ c],
\end{gathered}
\]
the two rank equalities of the statement describe the same condition. For a positive semidefinite
$\Gamma$ the elementary symmetric functions of its eigenvalues satisfy
$e_j(\Gamma)>0\iff\operatorname{rank}\Gamma\ge j$ \cite[Ch.~4]{hj13}, so the displayed disjunction
holds exactly when $t=\operatorname{rank}(MM^{\mathsf T})$, and
$e_{t+1}(\Gamma')=0$ together with $\Gamma'\succeq\Gamma$ forces
$\operatorname{rank}(MM^{\mathsf T}+cc^{\mathsf T})=t$.
\end{proof}

\paragraph*{Spectral reduction}\label{sec:spectral}
Since $B$ has full row rank,
\[
\operatorname{rank}P(\lambda)
=\operatorname{rank}\bigl((\lambda X-G)B\bigr)
=\operatorname{rank}(\lambda X-G).
\]
Thus \eqref{eq:5b} is equivalent to
\[
\operatorname{rank}(\lambda X-G)=\omega,\qquad\forall\lambda\in\mathbb C .
\]

\subsection{The spectral certificate}\label{sec:spectrum}

\begin{lem}[B\'ezout degree bound]\label{lem:bezout}\label{sec:bezout}
	Let $\omega\ge1$, and let $f_1,\ldots,f_s\in\mathbb R[\lambda]$ have
	degree at most $\omega$ and no common complex root. Then there exist
	real polynomials $g_i$ with $\deg g_i<\omega$ such that
	\[
	\sum_i g_i f_i=1 .
	\]
\end{lem}

\begin{proof}
	Choose $a\in\mathbb R^s$ such that
	\[
	g=\sum_i a_if_i\neq0,
	\]
	which is possible because the $f_i$ have no common complex root. Let
	$\alpha_1,\dots,\alpha_m$ be the distinct complex roots of $g$, and for
	$b\in\mathbb R^s$ put
	\[
	h=\sum_i b_if_i,\qquad
	h(\alpha_j)=\ell_j(b)+i\,m_j(b),
	\]
	where $\ell_j$ and $m_j$ are real linear forms. Some $f_i(\alpha_j)$ is
	nonzero, since otherwise $\alpha_j$ would be a common root, so the pair
	$(\ell_j,m_j)$ does not vanish identically and
	\[
	V_j:=\{b\in\mathbb R^s:\ h(\alpha_j)=0\}
	=\ker\ell_j\cap\ker m_j\subsetneq\mathbb R^s .
	\]
	A finite union of proper subspaces cannot cover $\mathbb R^s$, hence
	\[
	\begin{gathered}
	\exists\,b\in\mathbb R^s:\quad b\notin\bigcup_{j=1}^{m}V_j,\\
	\text{that is,}\quad h(\alpha_j)\neq0\quad (j=1,\dots,m).
	\end{gathered}
	\]
	Then $g$ and $h$ have no common root and are therefore coprime, so the two-polynomial
	B\'ezout identity with degree bounds \cite[Ch.~3]{vzgg13} yields
	$u,v\in\mathbb R[\lambda]$ with
	\[
	\begin{gathered}
	ug+vh=1,\\
	\deg u<\deg h\le\omega,\qquad \deg v<\deg g\le\omega .
	\end{gathered}
	\]
	Expanding $g,h$ in the $f_i$ and setting $g_i=ua_i+vb_i$ gives
	\[
	\sum_i g_i f_i=1,\qquad \deg g_i<\omega .
	\]
\end{proof}

\begin{lem}[Window]\label{lem:window}
	Let $f_1,\dots,f_s\in\mathbb R[\lambda]$ have degree at most $\omega$,
	let $c_{i,a}$ be the coefficient of $\lambda^a$ (with $c_{i,a}=0$ for
	$a<0$ or $a>\omega$), and let $Z$ be the matrix with rows $(i,j)$,
	$0\le j<\omega$, columns $t=0,\dots,2\omega-1$, and
	$Z_{(i,j),t}=c_{i,t-j}$. For a polynomial $p$ of degree at most
	$2\omega-1$, let $\operatorname{coef}(p)\in\mathbb R^{2\omega}$ be its
	coefficient vector in the basis $1,\lambda,\dots,\lambda^{2\omega-1}$.
	Then
	\[
	\operatorname{row}Z=
	\Bigl\{
	\operatorname{coef}\Bigl(\textstyle\sum_i g_if_i\Bigr)
	:\ \deg g_i<\omega
	\Bigr\}.
	\]
\end{lem}

\begin{proof}
	Row $(i,j)$ of $Z$ is $\operatorname{coef}(\lambda^jf_i)$, because
	\[
	\lambda^jf_i=\sum_t c_{i,t-j}\lambda^t .
	\]
	For weights $y_{i,j}$, the linearity of $\operatorname{coef}$ on
	polynomials of degree at most $2\omega-1$ gives
	\[
	\begin{aligned}
	\sum_{i,j}y_{i,j}\operatorname{coef}(\lambda^jf_i)
	&=\operatorname{coef}\Bigl(\sum_{i,j}y_{i,j}\lambda^jf_i\Bigr)\\
	&=\operatorname{coef}\Bigl(\sum_i g_if_i\Bigr),
	\qquad g_i=\sum_j y_{i,j}\lambda^j,
	\end{aligned}
	\]
	with $\deg g_i<\omega$; conversely, every family $g_i$ with
	$\deg g_i<\omega$ arises from its own coefficients $y_{i,j}$. Since
	$\deg(\lambda^jf_i)\le\omega+j\le2\omega-1$, all coefficients fit in
	the $2\omega$ columns, and the convention $c_{i,a}=0$ outside
	$0\le a\le\omega$ makes the truncation automatic.
\end{proof}

\begin{prop}[Gram certificate]\label{prop:gram}\label{sec:gram}
	Let $\omega\ge1$ and let $f_i$ be the maximal minors of $P(\lambda)$.
	Form a matrix $Z$ whose rows are indexed by $(i,j)$, $0\le j<\omega$,
	and whose columns are indexed by $t=0,\ldots,2\omega-1$, with
	\[
	Z_{(i,j),t}=c_{i,t-j},
	\]
	where $c_{i,a}$ is the coefficient of $\lambda^a$ in $f_i$, and let
	$e_0=(1,0,\ldots,0)^{\mathsf T}\in\mathbb R^{2\omega}$. Then
	\[
	\eqref{eq:5b}
	\iff e_0^{\mathsf T}\in\operatorname{row}Z
	\iff\exists x\in\mathbb R^{2\omega}:\ Z^{\mathsf T}Zx=e_0 .
	\]
\end{prop}

\begin{proof}
	By construction, $\operatorname{coef}(\lambda^jf_i)$ is row $(i,j)$ of
	$Z$, so Lemma~\ref{lem:window} gives
	\[
	\operatorname{row}Z=
	\Bigl\{
	\operatorname{coef}\Bigl(\textstyle\sum_i g_if_i\Bigr)
	:\ \deg g_i<\omega
	\Bigr\},
	\]
	and with $e_0^{\mathsf T}=\operatorname{coef}(1)$,
	\[
	e_0^{\mathsf T}\in\operatorname{row}Z
	\iff
	\exists\,g_i\in\mathbb R[\lambda],\ \deg g_i<\omega:\
	\sum_i g_if_i=1 .
	\]
	The right-hand side is the B\'ezout representation of
	Lemma~\ref{lem:bezout}, which exists if and only if the $f_i$ have no
	common complex root, that is, if and only if \eqref{eq:5b} holds. Thus
	$\operatorname{row}Z$ is the degree-windowed part of the ideal generated
	by the $f_i$, not the whole ideal, and the window $\deg g_i<\omega$ is
	exactly the one supplied by Lemma~\ref{lem:bezout}; no bound on the
	degree of an arbitrary relation is needed. Finally, for real $Z$,
	\[
	\operatorname{row}Z=\operatorname{im}Z^{\mathsf T}
	=\operatorname{im}(Z^{\mathsf T}Z),
	\]
	the last equality because $\ker Z=\ker(Z^{\mathsf T}Z)$; hence
	$e_0^{\mathsf T}\in\operatorname{row}Z$ if and only if
	$Z^{\mathsf T}Zx=e_0$ has a real solution.
\end{proof}

\paragraph*{Cauchy--Binet compression}\label{sec:cb}
Work in the coordinates of Remark~\ref{rem:YS}, in which the first $\ell$ vectors are a basis of
$S=\operatorname{row}(AN)$ and the remaining $k-\ell$ ones span a complement. For a feasible $W$ the
two projections coincide and lie in $S$ (Lemma~\ref{lem:struct}), so every row of $H=LN$ and of
$H_A=LAN$ belongs to $S$; hence in these coordinates the last $k-\ell$ columns of $P(\lambda)=\lambda
H-H_A$ vanish, and a nonzero minor uses only the $\ell$ columns of $S$. Accordingly define
\[
K(\lambda,\mu)=\sum_{|J|=\omega}
\det\left(P(\lambda)_{:,J}^{\mathsf T}P(\mu)_{:,J}\right),
\]
where $J\subseteq[\ell]$; only balanced column sets matter, since a column set meeting the
complement of $S$ gives a zero determinant. For each fixed column set $J$, the Cauchy--Binet identity
\cite[Ch.~0]{hj13} applied
to the $\omega\times q$ matrix $P(\lambda)_{:,J}^{\mathsf T}$ and the
$q\times\omega$ matrix $P(\mu)_{:,J}$ gives
\[
\det\!\left(P(\lambda)_{:,J}^{\mathsf T}P(\mu)_{:,J}\right)
=\sum_{\substack{I\subseteq[q]\\ |I|=\omega}}
f_{I,J}(\lambda)f_{I,J}(\mu),
\]
where $f_{I,J}$ is the minor on rows $I$ and columns $J$. Summing over $J$,
\[
K(\lambda,\mu)=\sum_{\substack{I\subseteq[q]\\ |I|=\omega}}
\sum_{\substack{J\subseteq[\ell]\\ |J|=\omega}}
f_{I,J}(\lambda)f_{I,J}(\mu),
\]
the Gram kernel of \textbf{all} maximal minors of $P$.

Three consequences are used below.
\begin{enumerate}[label=(\roman*)]
	\item\label{item:cost} \emph{Cost.} Computing $K$ needs only
	$\binom{\ell}{\omega}\le2^{\ell}$ determinants of order $\omega$, one per
	column set $J$ of the $\ell$-dimensional space $S$ of \eqref{eq:ell}; the
	$\binom q\omega\binom k\omega$ individual minors are never enumerated,
	and no hidden enumeration replaces them.
	\item\label{item:coef} \emph{Coefficient identity.} Writing $S_{a,b}$
	for the coefficient of $\lambda^a\mu^b$ in $K$ and $c_{i,a}$ for the
	coefficient of $\lambda^a$ in $f_i$, one has
	$S_{a,b}=\sum_i c_{i,a}c_{i,b}$. Hence, with the row indexing of
	Lemma~\ref{lem:window},
	\[
	(Z^{\mathsf T}Z)_{t,t'}
	=\sum_{j=0}^{\omega-1}\sum_i c_{i,t-j}c_{i,t'-j}
	=\sum_{j=0}^{\omega-1}S_{t-j,t'-j}.
	\]
	\item\label{item:roots} \emph{Common roots.} Since the coefficients of
	$P$ are real, specializing at $(\lambda,\mu)=(\alpha,\bar\alpha)$ gives
	\[
	K(\alpha,\bar\alpha)
	=\sum_{I,J}\left|f_{I,J}(\alpha)\right|^2,
	\]
	so $K(\alpha,\bar\alpha)=0$ if and only if every maximal minor of $P$
	vanishes at $\alpha$. This identity shows that the coefficient matrix of
	$K$ encodes precisely the common-root condition; the algorithm below
	nevertheless decides that condition at the coefficient level, without
	evaluating $K$ at complex points.
\end{enumerate}

\begin{remark}[Algorithmic reading]\label{rem:algo}
	The compression is computed in four steps:
	(i) for each column set $J\subseteq[k]$, $|J|=\omega$, form the
	$\omega\times\omega$ pencil $P(\lambda)_{:,J}^{\mathsf T}P(\mu)_{:,J}$;
	(ii) compute its bivariate coefficients, of degree at most $\omega$ in
	each variable;
	(iii) sum the coefficients over $J$;
	(iv) recover $Z^{\mathsf T}Z$ from the shifted identity of
	item~\ref{item:coef}.
	The identity is classical (an exterior-power Gram identity; see, e.g.,
	\cite{emi99} for the B\'ezoutian/elimination background). What the
	construction adds is the combination with the parameter-sized
	certificate: by the cost bound of item~\ref{item:cost}, the certificate
	matrix stays parameter-sized when $q$ is large and $k$ is bounded.
\end{remark}

\section{The fixed-order algorithm}\label{sec:fixedq}

This section decides $\mathrm{DEC}(q)$ at a prescribed order $q$. The procedure
is Algorithm~\ref{alg:fixed-q}: in each pivot chart it quantifies only over the
entries of $X$, $G$ and the certificate vector, that is over
$2q\omega+2\omega=O(q^2)$ real variables, whereas the chart coordinates of the
order-by-order procedure of \cite{fer26b} number $d(m_L-d)=(n-p-q)(q-r)$ and
still grow linearly in $n$ when $q$ is fixed; the compression of
Section~\ref{sec:compressed} is what removes that dependence.

\begin{algorithm}[!t]
\caption{Deciding $\mathrm{DEC}(q)$ for a fixed order $q$}
\label{alg:fixed-q}
\begin{algorithmic}[1]
\Require $A\in\mathbb Q^{n\times n}$, $C\in\mathbb Q^{p\times n}$,
$L_0\in\mathbb Q^{r_0\times n}$, integer $q\ge r_0$.
\Ensure \textsc{feasible} or \textsc{infeasible}.
\State $N\gets$ basis of $\ker[C;CA]$;\quad $k\gets$ number of columns of $N$
\State $h^0\gets L_0N$;\quad $g^0\gets L_0AN$;\quad
$D\gets$ rational matrix with $\ker D=\operatorname{row}[N\ AN]$
\For{$\omega=0$ \textbf{to} $\min(q,k)$}
  \For{each $J\subseteq[k]$ with $|J|=\omega$}
    \State form $\Phi_{\omega,J}(X,G,x)$ over $X,G\in\mathbb R^{q\times\omega}$,
           $x\in\mathbb R^{2\omega}$ with
    \Statex \quad (i) $e_\omega(X^{\mathsf T}X)>0$ (the left inverse $U$ is
           not quantified);
    \Statex \quad (ii) the rank-consistency formula of
           Lemma~\ref{lem:rankcons} for $M(X,G)b=c(X,G)$, $B_{:,J}=I_\omega$;
    \Statex \quad (iii) $X_iB=h_i^0$ and $G_iB=g_i^0$ for $i\le r_0$;
    \Statex \quad (iv) $D[(X_iB)^{\mathsf T};(G_iB)^{\mathsf T}]=0$ for
           $i>r_0$;
    \Statex \quad (v) $Z^{\mathsf T}Zx=e_0$, with $Z$ from the maximal minors
           of $\lambda X-G$ (Proposition~\ref{prop:gram}).
    \If{$\Phi_{\omega,J}$ is real-feasible}
      \State \Return \textsc{feasible}
    \EndIf
  \EndFor
\EndFor
\State \Return \textsc{infeasible}
\end{algorithmic}
\end{algorithm}

\begin{thm}\label{thm:fixedq}
For every fixed $q$, $\mathrm{DEC}$ is decidable in deterministic polynomial
time. More precisely, Algorithm~\ref{alg:fixed-q} runs in
$\mathcal L^{O(q^2)}$ bit operations.
\end{thm}

\begin{proof}
\emph{(i) Equivalence.} Let $L=[L_0;R]$ be feasible, $H=LN$, $H_A=LAN$,
$\omega=\operatorname{rank}H$, and let $J$ be a pivot set with
\[
H=XB,\qquad B_{:,J}=I_\omega .
\]
By Lemma~\ref{lem:compression}(ii) there is $F$ with $H_A=FH$; put $G=FX$.
Then
\[
X_iB=H_i=h_i^0,\qquad G_iB=(H_A)_i=g_i^0\qquad(i\le r_0),
\]
and, since $\ker D=\mathcal T=\operatorname{row}[N\ AN]$,
\[
D[(X_iB)^{\mathsf T};(G_iB)^{\mathsf T}]=0\qquad(i>r_0).
\]
Moreover $\operatorname{rank}X=\omega$ gives $e_\omega(X^{\mathsf T}X)>0$,
Lemma~\ref{lem:rankcons} gives the rank-consistency condition for
$B_{:,J}=I_\omega$, and Proposition~\ref{prop:gram} turns \eqref{eq:5b} into
$Z^{\mathsf T}Zx=e_0$. Hence $\Phi_{\omega,J}$ holds.

Conversely, suppose $\Phi_{\omega,J}$ holds. Lemma~\ref{lem:rankcons} yields
$B$ with $B_{:,J}=I_\omega$; put
\[
H=XB,\qquad H_A=GB .
\]
The fixed-row equations give $H_i=h_i^0$ and $(H_A)_i=g_i^0$ for $i\le r_0$,
and the free-row equations place each remaining pair in $\mathcal T$, so there
are $r_i$ with
\[
(H_i,(H_A)_i)=r_i[N\ AN]\qquad(i>r_0).
\]
With $L=[L_0;R]$, $R$ the matrix with rows $r_i$, this reads $H=LN$ and
$H_A=LAN$. Since $e_\omega(X^{\mathsf T}X)>0$, $X$ has a left inverse $U$;
hence
\[
\begin{aligned}
G&=(GU)X,\\
H_A&=GB=(GU)XB=FH,\qquad F=GU,
\end{aligned}
\]
so $\operatorname{row}H_A\subseteq\operatorname{row}H$, which is
\eqref{eq:5a} by Lemma~\ref{lem:compression}(i), and $Z^{\mathsf T}Zx=e_0$ is
\eqref{eq:5b} by Proposition~\ref{prop:gram}. Therefore the disjunction
$\bigvee_{\omega,J}\Phi_{\omega,J}$ holds if and only if a feasible $L$
exists, which is the test performed by Algorithm~\ref{alg:fixed-q}.

\emph{(ii) Count.} With $\omega\le k$ and $N=2q\omega+2\omega$,
\begin{align}
N &= 2q\omega+2\omega\le 2q^2+2q,\\
s &= O(k)\cdot O(k)+O(1)=O(k^2),\\
D &= \max\{O(k),\,2\omega+1\}=O(k)\le O(n),\\
B &= \operatorname{poly}(n,\mathcal L),
\end{align}
where the two factors of $s$ count the rank-consistency branches
$t\le\min\{\text{rows},\text{cols}\}(M)$, at most $O(k)$ of them, and the
$O(k)$ polynomials of degree $O(k)$ carried by each branch, while $2\omega+1$
is the degree of the Gram and certificate constraints; the last line is the
coefficient-size lemma of Section~\ref{sec:fpt}. Renegar's Part II bound
$B(\log B)(\log\log B)(sD)^{O(N)}$ therefore gives, per chart,
\[
B(\log B)(\log\log B)(sD)^{O(N)}
=\operatorname{poly}(n)^{O(q^2)}
=\mathcal L^{O(q^2)},
\]
because $sD=\operatorname{poly}(n)$ and $N\le 2q^2+2q$ with $q$ fixed. The
charts number
\[
N_{\mathrm{ch}}=(q+1)k^{q}\le 2^{q}k^{q}\le\mathcal L^{2q},
\]
since $k\le n\le\mathcal L$ and $q+1\le 2^{q}\le\mathcal L^{q}$ for $q\ge1$
and $\mathcal L\ge2$. Hence
\[
N_{\mathrm{ch}}\cdot\mathcal L^{O(q^2)}
\le\mathcal L^{2q}\cdot\mathcal L^{O(q^2)}
=\mathcal L^{O(q^2)} .
\]
For fixed $q$ the exponent $O(q^2)$ is a constant, so the construction size is
$k^{O(q^2)}\operatorname{poly}(\mathcal L)=\mathcal L^{O(q^2)}$ and Renegar's
bound gives the same order. The rank index $t$ is a discrete branch, not a
degree: it is counted in $s$, not folded into $D$. The degenerate cases $k=0$
and $q=r_0$ are treated in Section~\ref{sec:deg}.
\end{proof}

Table~\ref{tab:4} collects the quantities entering the bound.

\begin{table*}[!t]
\centering
\caption{Quantities entering the fixed-order bound.}\label{tab:4}
\small
\begin{tabular}{p{0.450\textwidth}p{0.450\textwidth}}
\toprule
Quantity & Bound \\
\midrule
real variables $N$ & $2q\omega+2\omega\le2q^2+2q$ \\
pivot charts & $(q+1)k^q$ \\
rank-consistency branches $t$ & $O(k)$, one per possible rank \\
number of polynomials $s$ (each counted once) & $O(k^2)$ \\
maximum degree $D$ & $O(k)$ \\
coefficient bit length $B$ & $\operatorname{poly}(n,\mathcal L)$ \\
construction size & $k^{O(q^2)}\operatorname{poly}(\mathcal L)$ \\
\bottomrule
\end{tabular}
\end{table*}

Three representation details are worth recording.

\begin{enumerate}
\item \emph{Strict inequalities.} A condition $f>0$ is written as
\[
f>0\iff\exists u:\ fu^2=1,
\]
which is equivalent over the reals. There are $O(1)$ strict inequalities per
branch, so this adds $O(1)$ variables and changes no order.
\item \emph{Disjunctions.} The rank-consistency disjunction is processed branch
by branch, with one real-feasibility query per $t$ and the running times
added:
\[
t\cdot\mathcal L^{O(q^2)}=O(k)\,\mathcal L^{O(q^2)}=\mathcal L^{O(q^2)} .
\]
This is why $t$ appears in $s$ and not in $D$.
\item \emph{Elementary symmetric polynomials.} $e_j(G)$ is represented by the
coefficients of $\det(zI+G)$, computed by fraction-free elimination on the
$k\times k$ polynomial matrix $zI+G$ in $O(k^3)$ polynomial operations of
degree $O(k)$. The expansion is never written out term by term.
\end{enumerate}

\subsection{Degenerate cases}\label{sec:deg}

\begin{enumerate}
\item \emph{$k=0$}, that is, $\ker[C;CA]=0$. Then $N$ has no columns, so
\[
H=LN=0,\qquad H_A=LAN=0\qquad\text{for every }L,
\]
and both \eqref{eq:5a} and \eqref{eq:5b} hold; in particular for $L=L_0$, so
$q_{\min}=r_0$.
\item \emph{$\omega=0$ with $k>0$.} Then $H=0$, and
\[
\operatorname{rank}(\lambda H-H_A)=0\ \ \forall\lambda\in\mathbb C
\iff H_A=0,
\]
which together with $H=0$ is exactly $\eqref{eq:5a}\wedge\eqref{eq:5b}$ for
that $L$. The branch $H=H_A=0$ of the encoding is therefore neither stronger
nor weaker than the condition.
\item \emph{$q<r_0$.} Rejected by definition, since $L=[L_0;R]$ has $q$ rows
and contains the $r_0$ rows of $L_0$.
\item \emph{$r=0$.} Decided by the exact pencil test of Section~\ref{sec:zero}.
\end{enumerate}

\section{The fixed-parameter algorithm}\label{sec:fpt}

By Proposition~\ref{prop:minorder} a feasible observer is a subspace $W$ with
$W_0\subseteq W\subseteq\mathcal T$, $\dim(W/W_0)\le r$, and the two conditions of
Lemma~\ref{lem:W-conditions}. Lemma~\ref{lem:struct} confines the feasible $W$ to
$\omega+1\le\dim W\le2\omega$ with $\omega\le\ell$, and Remark~\ref{rem:YS} puts every object in
$\mathbb R^\ell$. The algorithm therefore quantifies over a chart of a Grassmannian of dimension
at most $2\ell$ and over the $2\omega$ certificate variables, and the only non-polynomial factor
is the parameter $\ell=\operatorname{rank}(AN)$, which by \eqref{eq:ell} does not involve $r$.

\begin{algorithm}[!t]
\caption{Deciding $\mathrm{DEC}(q)$ in time $f(\ell)\operatorname{poly}(\mathcal L)$}
\label{alg:fpt}
\begin{algorithmic}[1]
\Require $A,C,L_0,q$ as above; $r=q-r_0$.
\Ensure \textsc{feasible} or \textsc{infeasible}.
\If{$r=0$}
  \State \Return the exact test of Proposition~\ref{prop:zero} on $L=L_0$
\EndIf
\If{$q<q_0$}
  \State \Return \textsc{infeasible}
  \Statex \quad \textup{(the algebraic lower bound \eqref{eq:range})}
\EndIf
\State $N\gets$ basis of $\ker[C;CA]$;\quad $M\gets[N\ AN]$;\quad
$S\gets\operatorname{row}(AN)$;\quad $\ell\gets\operatorname{rank}(AN)$
\State $W_0\gets\operatorname{row}(L_0M)$;\quad
$\mathcal V\gets\mathcal T\cap(S\times S)$
\If{$W_0\not\subseteq\mathcal V$}
  \State \Return \textsc{infeasible}
  \Statex \quad \textup{(necessary: every feasible $W$ satisfies $W_0\subseteq W\subseteq S\times S$)}
\EndIf
\If{$W_0=0$}
  \State \Return \textsc{feasible}
  \Statex \quad \textup{(the $\omega=0$ branch: $W=0$ is realized by $R=0$)}
\EndIf
\State $s_0\gets\dim W_0$;\quad $m'\gets\dim\mathcal V-s_0$
\For{$t=0$ \textbf{to} $\min(r,2\ell)$}
  \For{each chart of $\mathrm{Gr}(t,m')$ with chart coordinates $Y\in\mathbb R^{t\times(m'-t)}$}
    \State $H(Y)\gets$ first block of $Z(Y)$;\quad $G(Y)\gets$ second block
    \For{$\omega=1$ \textbf{to} $\ell$}
      \State form $\Psi_{t,\omega}(Y,x)$ with
      \Statex \quad (i) $\sigma_\omega(HH^{\mathsf T})>0$ and
               $\sigma_{\omega+1}\bigl([H;G][H;G]^{\mathsf T}\bigr)=0$;
      \Statex \quad (ii) $B(Y)^{\mathsf T}B(Y)x=\mathbf e_1$, with $B^{\mathsf T}B$ computed
               from the kernel $K(\lambda,\mu)$ by Cauchy--Binet (Lemma~\ref{lem:window}) over the
               $\ell$ columns of $S$.
      \If{$\Psi_{t,\omega}$ is real-feasible}
        \State \Return \textsc{feasible}
      \EndIf
    \EndFor
  \EndFor
\EndFor
\State \Return \textsc{infeasible}
\end{algorithmic}
\end{algorithm}

\begin{thm}\label{thm:fpt}
For rational input and every $r\ge0$, $\mathrm{DEC}$ is decidable in
\begin{equation}\label{eq:fpt-l}
\exp\bigl(O(\ell^2\log(\ell+2))\bigr)\operatorname{poly}(\mathcal L)
\end{equation}
bit operations, where $\ell=\operatorname{rank}(AN)\le k$ does not depend on the number $r$ of
free observer rows. For $r=0$ the bound sharpens to $\operatorname{poly}(\mathcal L)$.
\end{thm}

\begin{proof}
The case $r=0$ is Proposition~\ref{prop:zero}. For $r\ge1$ the argument has six steps.

\emph{(i) Reduction to a subspace.} By Lemma~\ref{lem:sandwich} a feasible free block exists if and
only if some feasible $W$ satisfies $W_0\subseteq W\subseteq\mathcal T$ and
$\dim(W/W_0)\le r$, and by Lemma~\ref{lem:monotone} the bound $r$ can be saturated; it is therefore
enough to range $t=\dim(W/W_0)$ over $0,\dots,\min(r,2\ell)$ and to test each $t$ once.

\emph{(ii) The inclusion test is necessary.} If $\pi_1(W)=\pi_2(W)=U\subseteq S$
(Lemma~\ref{lem:struct}), then $W\subseteq U\times U\subseteq S\times S$, and $W\subseteq\mathcal T$
by Lemma~\ref{lem:sandwich}; hence $W_0\subseteq W\subseteq\mathcal V=\mathcal T\cap(S\times S)$.
The test at the head of the algorithm therefore rejects no feasible instance, and it is what makes
$m'=\dim\mathcal V-s_0$ a nonnegative number.

\emph{(iii) The bound $\omega\le\ell$ and $t\le2\ell$.} Since $U\subseteq S$ we have
$\omega=\dim U\le\ell$, and by Lemma~\ref{lem:struct} $\dim W\le2\omega\le2\ell$, so
$t\le2\ell$.

\emph{(iv) The charts exhaust the candidates.} Every $W$ with
$W_0\subseteq W\subseteq\mathcal V$ maps to a subspace $W/W_0$ of $\mathcal V/W_0$ of dimension
$t$, hence lies in one of the charts of $\mathrm{Gr}(t,m')$ that the algorithm enumerates.

\emph{(v) The formula characterizes the two conditions.} Within a chart, $\Psi_{t,\omega}$ is
satisfiable exactly when the row space of $Z(Y)$ satisfies \eqref{eq:5a} with
$\operatorname{rank}H=\omega$ and \eqref{eq:5b}, by Lemma~\ref{lem:compression} and
Lemma~\ref{lem:window}.

Fix $t$. By Lemma~\ref{lem:struct} and Remark~\ref{rem:YS} a feasible $W$ of quotient dimension
$t$ is a point of the Grassmannian $\mathrm{Gr}(t,m')$ with $m'\le2\ell-s_0$, written in the chart
coordinates $Y\in\mathbb R^{t\times(m'-t)}$ as the row space of $Z(Y)=[E_0;[I_t\mid Y]E_Q]$, which
contains $W_0$ by construction. The chart count is
\begin{equation}\label{eq:charts}
\sum_{t}\binom{m'}{t}=2^{m'}\le4^{\ell}.
\end{equation}

Within a chart and for a fixed $\omega$, the quantified variables are $Y$ and the certificate
$x\in\mathbb R^{2\omega}$, and the data of the real-feasibility problem are
\begin{equation}\label{eq:nvar}
\begin{gathered}
N=t(m'-t)+2\omega\le\tfrac{m'^2}{4}+2\omega\le\ell^2+2\ell,\\
s=2\omega+O(1)\le2\ell+O(1),\qquad D=O(\ell),\\
B=\operatorname{poly}(\mathcal L),
\end{gathered}
\end{equation}
where the discrete choices of $t$, of the chart and of $\omega$ are finite disjunctions rather
than real variables, and $B$ is the coefficient bit length of Lemma~\ref{lem:coeff}. The
polynomials are those of Lemma~\ref{lem:window}: $\sigma_\omega(HH^{\mathsf T})>0$ and
$\sigma_{\omega+1}([H;G][H;G]^{\mathsf T})=0$ encode \eqref{eq:5a} together with
$\operatorname{rank}H=\omega$ (Lemma~\ref{lem:compression}), and $B(Y)^{\mathsf T}B(Y)x=\mathbf
e_1$ encodes \eqref{eq:5b}. Substituting \eqref{eq:nvar} into \eqref{eq:rengar} gives, per
chart,
\begin{equation*}
\begin{gathered}
T\le B(\log B)(\log\log B)(sD)^{O(N)}\\
=\operatorname{poly}(\mathcal L)\exp\bigl(O(\ell^2\log(\ell+2))\bigr),
\end{gathered}
\end{equation*}
and multiplying by the chart count \eqref{eq:charts} absorbs $\exp(O(\ell))$ into the same
exponential. The total is \eqref{eq:fpt-l}.

\emph{(vi) Conclusion.} By (iv) every feasible $W$ is reached, and by (v) the formula accepts exactly
those $W$ that satisfy the two conditions, so the algorithm accepts if and only if a feasible observer
exists.
\end{proof}

\begin{cor}\label{cor:arith}
The same algorithm uses $\exp\bigl(O(\ell^2\log(\ell+2))\bigr)\operatorname{poly}(\mathcal L)$
arithmetic operations.
\end{cor}

\subsection{Minimum order and construction}\label{sec:optfind}

\begin{cor}[$\mathrm{OPT}$]\label{cor:opt}
$\mathrm{OPT}$ is decidable in $\exp\bigl(O(\ell^2\log(\ell+2))\bigr)\operatorname{poly}(\mathcal L)$
bit operations.
\end{cor}

Algorithm~\ref{alg:opt} makes the search explicit: $\mathrm{DEC}$ is used as a subroutine, and the
monotonicity of Lemma~\ref{lem:monotone} is what makes a binary search correct.

\begin{algorithm}[!t]
\caption{Minimum order $\mathrm{OPT}$, with $\mathrm{DEC}$ as a subroutine}
\label{alg:opt}
\begin{algorithmic}[1]
\Require $A,C,L_0,Q$; \ $\ell=\operatorname{rank}(AN)$; \ $r_0=\operatorname{rows}(L_0)$;
$q_0$ as in \eqref{eq:range}.
\Ensure $q_{\min}$, or $+\infty$ if no order up to $Q$ is feasible.
\State $\mathit{lo}\gets\max(0,\,q_0-r_0)$
\State $\mathit{hi}\gets\min\bigl(Q-r_0,\ n_0-p-r_0,\ 2\ell\bigr)$
\If{$\mathit{lo}>\mathit{hi}$ \textbf{or} $\mathrm{DEC}(r_0+\mathit{hi})$ returns \textsc{infeasible}}
  \State \Return $+\infty$
\EndIf
\While{$\mathit{lo}<\mathit{hi}$}
  \State $\mathit{mid}\gets\lfloor(\mathit{lo}+\mathit{hi})/2\rfloor$
  \If{$\mathrm{DEC}(r_0+\mathit{mid})$ returns \textsc{feasible}}
    \State $\mathit{hi}\gets\mathit{mid}$
  \Else
    \State $\mathit{lo}\gets\mathit{mid}+1$
  \EndIf
\EndWhile
\State \Return $r_0+\mathit{lo}$
\end{algorithmic}
\end{algorithm}

\begin{proof}
By Proposition~\ref{prop:minorder} $q_{\min}=r_0+t_{\min}$, and by \eqref{eq:range} and the same
proposition $t_{\min}$ lies in $[\max(0,q_0-r_0),\,\min(Q-r_0,\,n_0-p-r_0,\,2\ell)]$; that is the
initial interval of Algorithm~\ref{alg:opt}, so the search is confined before it starts. By
Lemma~\ref{lem:monotone} the predicate ``$\mathrm{DEC}(r_0+t)$ holds'' is monotone in $t$, so the
algorithm maintains the invariant that the least feasible $t$ lies in $[\mathit{lo},\mathit{hi}]$,
halves the interval at each step, and returns $t_{\min}$ after $O(\log\ell)$ calls to the decision
procedure of Theorem~\ref{thm:fpt}, the logarithmic factor being absorbed by \eqref{eq:fpt-l}. The
test at the head is what distinguishes ``no feasible order up to $Q$'' from a finite answer, and it
is correct because the predicate is upward closed.
\end{proof}

\begin{remark}[Room for improvement]\label{rem:refine}
The complexity bound does not depend on how tightly the window is chosen, so
Algorithm~\ref{alg:opt} uses only the published lower bound $q_0$ and the structural bound $2\ell$.
Two cheap refinements are available and are left out on purpose, to keep the algorithm free of
instance-specific preprocessing. First, \eqref{eq:qmin-lb} gives a lower bound that is computed in
polynomial time from $U_0$ and $V_0$ and can be strictly larger than $q_0$; on the instance of
Appendix~\ref{app:example} it equals $4$ while $q_0=3$. Replacing $\max(0,q_0-r_0)$ by the larger of
the two bounds in line~2 shrinks the window without affecting correctness. Second, the structural
bound \eqref{eq:struct} also confines $\omega$ on each chart, which prunes branches of the inner loop
in Algorithm~\ref{alg:fpt}. Neither refinement changes any asymptotic statement, and neither is
needed for the results above.
\end{remark}

\begin{cor}[$\mathrm{FIND}$, arithmetic model]\label{cor:find}
If the instance is feasible, a witness $R$ can be computed in
$\exp\bigl(O(\ell^2\log(\ell+2))\bigr)\operatorname{poly}(\mathcal L)$ arithmetic operations, as
exact real algebraic numbers whose number and degree are $f(\ell)$.
\end{cor}

Procedure~\ref{proc:find} states the construction; step~1 is the only step that is not elementary,
and it is a citation rather than a description.

\paragraph*{Procedure FIND (arithmetic model)}\label{proc:find}
\begin{enumerate}
\item Solve the system $\Psi_{t,\omega}$ of Section~\ref{sec:fpt} by the sampling-point algorithm of
\cite[Ch.~13]{bpr06}, for the least $t$ accepted by $\mathrm{DEC}$ and for some $\omega$ that
succeeds. This returns $Y$ and $x$ as exact real algebraic numbers, with the arithmetic bound of
\cite[Ch.~13]{bpr06}. We neither describe that algorithm nor bound its bit cost
(Remark~\ref{rem:find-bit}).
\item Form $W=\operatorname{row}Z(Y)$, choose a basis $w_1,\dots,w_t$ of $W/W_0$, and solve the
rational linear systems $u_iM=w_i$ for $i=1,\dots,t$; this costs $\operatorname{poly}(\mathcal L)$ by
Lemma~\ref{lem:rationalbasis}.
\item Return $R=[u_1;\dots;u_t;0]$ padded with $Q-r_0-t$ zero rows.
\end{enumerate}

\begin{proof}[Proof of Corollary~\ref{cor:find}]
Step~1 is the sampling-point bound of \cite[Ch.~13]{bpr06}: a nonempty real algebraic set defined by
$s$ polynomials in $N$ variables of degree at most $D$ and coefficient bit length $B$ has a sample
point whose coordinates are algebraic numbers of degree $D^{O(N)}$ with isolating intervals of bit
length $B D^{O(N)}$; with \eqref{eq:nvar} this is $f(\ell)\operatorname{poly}(\mathcal L)$ arithmetic
operations and $f(\ell)\operatorname{poly}(\mathcal L)$ output bits. Steps~2 and~3 are
Lemma~\ref{lem:sandwich}: a basis of $W/W_0$ lifts to free rows by solving $uM=w$, and padding with
zero rows realizes $W$ with exactly $t$ free rows. Since $W$ satisfies the two conditions
(Lemma~\ref{lem:W-conditions}), the resulting $L=[L_0;R]$ does too, and its order is $r_0+t$.
\end{proof}

\begin{remark}[What is not claimed]\label{rem:find-bit}
Corollary~\ref{cor:find} is an arithmetic bound on the construction together with a bound on the
output length. A bound in the bit model for producing the witness is not asserted: the ordinary
analysis of \cite{bpr06} gives arithmetic operations only, and their Chapter~8 states explicitly
that bit-length estimates do not follow from the arithmetic analysis. We therefore state
$\mathrm{FIND}$ in the arithmetic model and leave its bit complexity open
(Section~\ref{sec:limits}).
\end{remark}

\subsection{The bit model and the coefficient bound}\label{sec:bitsize}

We use the standard bit model: the input has rational entries of total bit length $\mathcal L$,
and $n,p,q\le\mathcal L$. Since only the row space of $C$ enters both conditions, we replace $C$ by
a rational basis of its row space when it has more than $n$ rows; this preserves $\operatorname{row}C$
and $\operatorname{row}[C;CA]$, hence $K_0$ and \eqref{eq:5a}--\eqref{eq:5b}, costs
$\operatorname{poly}(\mathcal L)$ bit operations, and leaves an output matrix with at most $n$ rows.
We therefore assume $p\le n$ in what follows. Renegar's Part II bit bound
\cite{ren92} has the form
\begin{equation}\label{eq:rengar}
B(\log B)(\log\log B)(sD)^{O(N)},
\end{equation}
for integer-coefficient input, where $N$ is the number of real variables, $s$
the number of polynomials, $D$ the maximum degree, and $B$ the maximum
coefficient bit length. The exponential factor depends only on $N$; the
coefficient-length contribution is polynomial in $\log B$. We apply it in the following
normalization, which we state once and use throughout. (i) All coefficients are integers: every system
is cleared of denominators beforehand, which changes neither the solution set nor the degree, and
inflates $B$ by the common denominator only. (ii) Every equation and every inequality is counted once
in $s$; there is no separate count for the two kinds. (iii) A strict inequality $g>0$ is replaced by
$g\ge\varepsilon$ with one fresh variable $\varepsilon$ constrained by $\varepsilon>0$, which adds at
most $s$ variables and leaves $s$ and $D$ as they are. (iv) The bound covers the decision only; the
construction of the system is accounted for separately, by Lemma~\ref{lem:rationalbasis} and
Lemma~\ref{lem:coeff}.

\begin{lem}[Rational bases and annihilators]\label{lem:rationalbasis}
Let $M\in\mathbb Q^{m\times n}$ have entries of bit length at most $\tau$, let $r=\operatorname{rank}M$
and put $\tau_M=O\bigl(r(\tau+\log n)\bigr)$. Then a basis $N$ of $\ker M$, a basis $R$ of
$\operatorname{row}M$, and a matrix $D$ with $\operatorname{row}D=\{x\in\mathbb Q^m:xM=0\}$ can be
computed in $\operatorname{poly}(m,n,\tau)$ bit operations, and the entries of $N$, $R$ and $D$ have
bit length at most $\tau_M$.
\end{lem}

\begin{proof}
Run fraction-free Gaussian elimination on $M$ and keep the pivot rows and columns. Every entry of the
resulting reduced row echelon form is a ratio of minors of $M$: by Cramer's rule the determinant in
the numerator is an $s\times s$ minor with $s\le r$, and by Hadamard's inequality
\begin{equation*}
\operatorname{bits}\det M_{I,J}\le O\bigl(s(\tau+\log s)\bigr)=O\bigl(r(\tau+\log n)\bigr).
\end{equation*}
The pivot columns give $R$, the free columns give $N$ after back-substitution, and the left null
space $D$ is obtained from the same elimination applied to $M^{\mathsf T}$; the common denominator
is a product of at most $r$ pivot minors, which adds only $O(r\log r)$ bits. The elimination uses
$O(mnr)$ arithmetic operations on numbers of the displayed bit length, so the total is
$\operatorname{poly}(m,n,\tau)$ bit operations.
\end{proof}

\begin{lem}[Coefficient sizes]\label{lem:coeff}
Let the input entries have bit length at most $\mathcal L$, and let
$n\le\mathcal L$. Then:
\emph{(i)} the kernel basis $N$, the matrix $D$, and the entries of
$H,H_A,G_H,G_S$ have bit length $\operatorname{poly}(n,\mathcal L)$;
\emph{(ii)} every maximal minor $f_i$ of $P(\lambda)=\lambda H-H_A$, and every
coefficient of $f_i$, has bit length $\operatorname{poly}(n,\mathcal L)$;
\emph{(iii)} the entries of $Z^{\mathsf T}Z$ and of $\sigma_j(G_H),\sigma_j(G_S)$ for
$j\le2\ell$ have bit length $f(\ell)\operatorname{poly}(\mathcal L)$.
\end{lem}

\begin{proof}
Let $m=n^2+pn+r_0n$ be the number of rational input entries. With $p\le n$ and $r_0\le q\le n$ this
is $m\le3n^2=O(n^2)$, and $m\le\mathcal L$ because each entry has bit length at least one. Multiplying
by the product of their denominators turns the input into an integer matrix $A_{\rm in}$ all of whose
entries are bounded by $B_0$ in absolute value, where
\begin{equation*}
m\le3n^2=O(n^2),\qquad B_0=2^{O(m\mathcal L)}=2^{O(n^2\mathcal L)} .
\end{equation*}
Hadamard's inequality gives, for a $t\times t$ minor $M_{I,J}$ of $A_{\rm in}$
with $t\le n$,
\begin{equation*}
\begin{aligned}
|\det M_{I,J}|&\le(\sqrt t\,B_0)^t,\\
\operatorname{bits}\det M_{I,J}
&=O\bigl(t(\log B_0+\log t)\bigr)=O(n^3\mathcal L),
\end{aligned}
\end{equation*}
which is $\operatorname{poly}(n,\mathcal L)$ because $n\le\mathcal L$.

\emph{(i)} Apply Lemma~\ref{lem:rationalbasis} to $[C;CA]$, whose entries have bit length at most
$\mathcal L$ and whose rank is $n-k$: the basis $N$ of $\ker[C;CA]$ has entries of bit length
$O\bigl((n-k)(\mathcal L+\log n)\bigr)$. Applying the lemma again to $[N\ AN]$, whose entries now
have that bit length, gives a row-space basis of $[N\ AN]$ and hence the matrix $D$, with entries of
bit length $O\bigl(d_0((n-k)(\mathcal L+\log n)+\log k)\bigr)=\operatorname{poly}(n,\mathcal L)$;
the two eliminations cost $\operatorname{poly}(\mathcal L)$ bit operations.
Each entry of $H=LN$ is a sum of $n$ products of an input entry with an entry of
$N$, and each entry of $G_H$ or $G_S$ is a sum of at most $q$ such products, so
passing from $N$ to $H$ and then to $G_H,G_S$ adds $O(\log n)$ and then
$O(\log q)$ bits:
\begin{equation*}
\begin{aligned}
\operatorname{bits}H_{ij}&\le\mathcal L+\operatorname{poly}(n,\mathcal L)+O(\log n),\\
\operatorname{bits}(G_H)_{ij},\ \operatorname{bits}(G_S)_{ij}
&\le2\operatorname{poly}(n,\mathcal L)+O(\log q).
\end{aligned}
\end{equation*}

\emph{(ii)} A maximal minor $f_i$ of $P(\lambda)=\lambda H-H_A$ is a determinant
of an $\omega\times\omega$ matrix whose entries are linear in those of $H,H_A$;
Hadamard applied to its coefficient matrix bounds each of its $\omega+1$
coefficients:
\begin{equation*}
\deg f_i\le\omega,\qquad
\operatorname{bits}c_{i,a}=\operatorname{poly}(n,\mathcal L).
\end{equation*}

\emph{(iii)} $Z^{\mathsf T}Z$ is a sum of at most
$\binom{\ell}{\omega}\le2^{\ell}$ terms, each a product of two coefficients of
the $f_i$, and the elementary symmetric functions of the $2\ell\times2\ell$ Gram
matrices are sums of at most $(2\ell)^{j}\le(2\ell)^{2\ell}$ principal minors. A
sum of $T$ terms of bit length $b$ has bit length $b+\log T$, hence
\begin{equation*}
\begin{aligned}
\operatorname{bits}(Z^{\mathsf T}Z)_{ij}
&\le2\operatorname{poly}(n,\mathcal L)+\ell,\\
\operatorname{bits}\sigma_j(G_H),\ \operatorname{bits}\sigma_j(G_S)
&\le\operatorname{poly}(n,\mathcal L)+2\ell\log(2\ell),
\end{aligned}
\end{equation*}
which is $f(\ell)\operatorname{poly}(\mathcal L)$.
\end{proof}

It remains to say why the per-chart bound already gives \eqref{eq:fpt-l}: the enumeration of
$t$, of the charts and of $\omega$ contributes the factors $2\ell+1$, $4^{\ell}$ and $\ell$, all
absorbed by the exponential; every chart and every $\omega$ costs one call to the decision procedure
of \eqref{eq:nvar}, and the rational data of each call have bit length
$f(\ell)\operatorname{poly}(\mathcal L)$ by Lemma~\ref{lem:coeff}. The cases $k=0$ and $\omega=0$
are treated directly in Section~\ref{sec:deg}, and $r=0$ in Section~\ref{sec:zero}. For fixed $q$
the sharper parameter is $q$ itself, and the bound $\mathcal L^{O(q^2)}$ of
Section~\ref{sec:fixedq} applies instead.

\subsection{The zero-free-row case: an exact polynomial-time pencil test}\label{sec:zero}

For $r=0$ only two conditions remain, and both are decidable without a
parameter-dependent search. \eqref{eq:5a} is a constant-rank test on
$H=L_0N$, $H_A=L_0AN$ (Lemma~\ref{lem:compression}(i)); \eqref{eq:5b} is the
constancy of the rank of the degree-one pencil $\lambda X-G$. We first isolate
the pencil statement.

\begin{lem}\label{lem:zero}
Let $X,G\in\mathbb Q^{q\times\omega}$ with $q\ge\omega\ge1$, and put
$Q(\lambda)=\lambda X-G$. The following are equivalent:
\emph{(i)} $\operatorname{rank}Q(\lambda)=\omega$ for every
$\lambda\in\mathbb C$;
\emph{(ii)} the maximal $\omega\times\omega$ minors of $Q$ have no common
complex root;
\emph{(iii)} the maximal minors do not all vanish and their gcd in
$\mathbb Q[\lambda]$ is a nonzero constant;
\emph{(iv)} there is $S(\lambda)\in\mathbb Q[\lambda]^{\omega\times q}$ with
$S(\lambda)Q(\lambda)=I_\omega$ and $\deg S\le2\omega-2$;
\emph{(v)} there exist matrices
$S_0,\dots,S_{2\omega-2}\in\mathbb Q^{\omega\times q}$ satisfying the rational
linear system
\begin{equation}
\tag{R0}\label{eq:R0}
\begin{aligned}
-S_0G&=I_\omega,\\
S_{j-1}X-S_jG&=0\quad(1\le j\le2\omega-2),\\
S_{2\omega-2}X&=0 .
\end{aligned}
\end{equation}
\end{lem}

\begin{proof}
\emph{(i)$\Leftrightarrow$(ii).} Since $Q(\lambda)$ is $q\times\omega$, its rank drops below
$\omega$ exactly when every $\omega\times\omega$ minor vanishes at that point, which is the
classical rank--minor criterion \cite[Ch.~0]{hj13}.

\emph{(ii)$\Leftrightarrow$(iii).} The maximal minors have rational coefficients and are not all
zero. Over the algebraically closed field $\mathbb C[\lambda]$, a common root of the minors is a
root of their gcd, and a nonconstant gcd has a root; hence the minors have no common root exactly
when their gcd is a nonzero constant \cite[Ch.~IV]{lan02}.

\emph{(iii)$\Rightarrow$(iv).} This is the adjugate/B\'ezout construction of
Appendix~\ref{app:encoding}. For each $\omega$-element row set
$I\subseteq[q]$, let $E_I\in\mathbb Q^{\omega\times q}$ be the row selection,
$Q_I=E_IQ$ and $f_I=\det Q_I$. Then
\begin{equation*}
\operatorname{adj}(Q_I)E_IQ=f_II_\omega,\qquad
\sum_Ig_If_I=1\ \text{ with }\ \deg g_I<\omega
\end{equation*}
by Lemma~\ref{lem:bezout}, so that
\begin{equation*}
S=\sum_Ig_I\operatorname{adj}(Q_I)E_I
\quad\text{satisfies}\quad
S(\lambda)Q(\lambda)=I_\omega ,
\end{equation*}
and $\deg\operatorname{adj}(Q_I)\le\omega-1$ gives
$\deg S\le(\omega-1)+(\omega-1)=2\omega-2$.

\emph{(iv)$\Rightarrow$(v).} Write $S=\sum_{j=0}^{2\omega-2}S_j\lambda^j$. Then
$S(\lambda)(\lambda X-G)=I_\omega$, and comparing the coefficients of
$\lambda^0$, $\lambda^k$ ($1\le k\le2\omega-2$) and $\lambda^{2\omega-1}$ gives
\begin{equation*}
-S_0G=I_\omega,\qquad S_{k-1}X-S_kG=0,\qquad S_{2\omega-2}X=0,
\end{equation*}
which is \eqref{eq:R0}.

\emph{(v)$\Rightarrow$(i).} Evaluating the polynomial identity
$S(\lambda)Q(\lambda)=I_\omega$ at any $\lambda_0\in\mathbb C$ exhibits a left
inverse of $Q(\lambda_0)$:
\begin{equation*}
S(\lambda_0)Q(\lambda_0)=I_\omega
\Longrightarrow\operatorname{rank}Q(\lambda_0)\ge\omega,
\end{equation*}
while $\operatorname{rank}Q(\lambda_0)\le\omega$ because $Q(\lambda_0)$ has
$\omega$ columns; hence $\operatorname{rank}Q(\lambda_0)=\omega$.
\end{proof}

The size of \eqref{eq:R0} is fixed by the degree bound $\deg S\le2\omega-2$.
With $d=2\omega-2$ the unknowns are the entries of $d+1=2\omega-1$ matrices of
size $\omega\times q$, and each of the $d+2=2\omega$ matrix equations contributes
$\omega^2$ scalar equations, so
\begin{equation*}
N_{\rm var}=(2\omega-1)q\omega=O(q\omega^2),\qquad
N_{\rm eq}=(2\omega)\omega^2=2\omega^3 .
\end{equation*}
All coefficients of \eqref{eq:R0} are entries of $X$, $G$, or $0,\pm1$, so
clearing the denominators multiplies the system by an integer of bit length
$O(\mathcal L_0)$, where $\mathcal L_0$ is the input length of $(X,G)$. Exact
fraction-free elimination on the resulting integer system of size
$M=N_{\rm var}+N_{\rm eq}=O(\omega^3+q\omega^2)$ then satisfies
\begin{equation*}
\begin{aligned}
\text{entry bits}&=O(\mathcal L_0),\\
\text{minor bits}&=O\bigl(M(\mathcal L_0+\log M)\bigr)\ \text{(Hadamard)},\\
\text{arithmetic operations}&=O(M^3),\\
\text{bit operations}&=O(M^3)\cdot O\bigl(M(\mathcal L_0+\log M)\bigr)^2\\
&=\operatorname{poly}(M,\mathcal L_0).
\end{aligned}
\end{equation*}
Hence consistency of \eqref{eq:R0} is decided in
$\operatorname{poly}(q,\omega,\mathcal L_0)$ bit operations.

\begin{remark}[Construction and algorithm]\label{rem:construction}
The adjugate construction in (iii)$\Rightarrow$(iv) of Lemma~\ref{lem:zero}
enumerates the $\binom q\omega$ maximal minors, which is exponential in $q$; it
is used only to prove the existence of a certificate of degree at most
$2\omega-2$. It applies to the rectangular pencil $Q$ because each $Q_I$ is
square, so that $\deg S\le2\omega-2$ holds for every $q\ge\omega$ with no square
assumption on $Q$. The algorithm never forms the sum: it searches for the
coefficients of $S$ directly in \eqref{eq:R0}, whose size is polynomial in
$q,\omega$. This is the step that turns a certificate-size statement into a
decision algorithm.
\end{remark}

\begin{remark}[Higher degree]\label{rem:degree}
The same argument applies to pencils of degree $d$: if the entries of $Q$ have
degree at most $d$, then every maximal minor has degree at most $d\omega$,
B\'ezout gives multipliers with $\deg g_I<d\omega$, and the adjugate contributes
at most $d(\omega-1)$,
\begin{equation*}
\deg S\le(d\omega-1)+d(\omega-1)=2d\omega-d-1,
\end{equation*}
which is $2\omega-2$ for $d=1$. Degree bounds for polynomial-matrix left
inverses are studied classically \cite{fab90}; the bound above is the elementary
adjugate--B\'ezout one, and only the case $d=1$ is used in this paper. The
general bound may be useful if the observer model is extended. For $d=0$ the
pencil is constant and the statement is the rank test of
Lemma~\ref{lem:compression}(iii).
\end{remark}

\begin{prop}\label{prop:zero}
Let $r=0$, so $L=L_0$. Feasibility of $\eqref{eq:5a}\wedge\eqref{eq:5b}$ is
decidable in $\operatorname{poly}(n,\mathcal L)$ bit operations.
\end{prop}

\begin{proof}
Compute a basis $N$ of $\ker[C;CA]$ and set $H=L_0N$, $H_A=L_0AN$.

\emph{Step 1.} Check $\operatorname{rank}[H_A;H]=\operatorname{rank}H$; if it
fails, reject, since this is \eqref{eq:5a}
(Lemma~\ref{lem:compression}(i)).

\emph{Step 2.} Put $\omega=\operatorname{rank}H$. If $\omega=0$, then $H=0$ and
\eqref{eq:5b} reads $\operatorname{rank}(\lambda H-H_A)=0$ for all $\lambda$,
i.e.\ $H_A=0$; accept exactly in that case.

\emph{Step 3.} Otherwise take a rank factorization $H=XB$ with
$X\in\mathbb Q^{q\times\omega}$ of full column rank and
$B\in\mathbb Q^{\omega\times k}$ of full row rank. Step 1 gives
$\operatorname{row}H_A\subseteq\operatorname{row}H$
(Lemma~\ref{lem:compression}(i)), hence $H_A=FH$ for some $F$; set $G=FX$, so
$H_A=GB$.

\emph{Step 4.} Run the test of Lemma~\ref{lem:zero} on $(X,G)$.

\emph{Correctness.} Since $B$ is of full row rank,
\begin{equation*}
\operatorname{rank}(\lambda H-H_A)
=\operatorname{rank}\bigl((\lambda X-G)B\bigr)
=\operatorname{rank}(\lambda X-G)
\end{equation*}
for every $\lambda$, and \eqref{eq:5b} is
$\operatorname{rank}(\lambda H-H_A)=\omega$ for all $\lambda$; thus
Lemma~\ref{lem:zero} decides \eqref{eq:5b}. Step 1 is \eqref{eq:5a}, and step 2
is the exact $\omega=0$ reading of the same identity. The zero-free-row decision
is therefore the conjunction of one constant-rank test and one pencil test: if
\eqref{eq:5a} fails, the instance is infeasible and the spectral question is not
asked.

\emph{Complexity.} Step 1 is one rank computation, step 3 a rank factorization
and an exact linear solve, and step 4 costs
$\operatorname{poly}(q,\omega,\mathcal L_0)\subseteq\operatorname{poly}(n,\mathcal L)$
by Lemma~\ref{lem:zero}. Every step is a rank computation or an exact linear
solve on a matrix of at most $2\omega+1$ rows and $q$ columns whose entries have
bit length $\operatorname{poly}(n,\mathcal L)$ by Lemma~\ref{lem:coeff}, and
Gaussian elimination over $\mathbb Q$ on such a matrix costs
\begin{equation*}
\operatorname{poly}(n,\mathcal L)\ \text{ bit operations};
\end{equation*}
the test returns the coefficients $S_0,\dots,S_{2\omega-2}$ explicitly.
\end{proof}

\begin{remark}[Kronecker canonical form]\label{rem:kronecker}
The pencil $Q$ has full column rank at every finite complex value if and only if
its Kronecker form (in the sense of \cite{bee88}) contains no right singular
block and no finite regular elementary-divisor block. Left singular blocks and
\emph{infinite} regular elementary divisors are permitted: they do not produce a
rank drop at any finite $\lambda$. For instance
\begin{equation*}
Q(\lambda)=1,\qquad Q(\lambda)=I_2,\qquad
Q(\lambda)=\begin{bmatrix}1&\lambda\\0&1\end{bmatrix}
\end{equation*}
all satisfy the condition, the first with an infinite elementary divisor of size
one and the last with determinant the constant $1$. Equivalently, in the
language of Appendix~\ref{app:encoding}, the condition is that the ideal
generated by the maximal minors is the unit ideal; only the \emph{finite}
spectrum is constrained.
\end{remark}

\begin{remark}[Witness]\label{rem:witness}
If \eqref{eq:R0} is consistent, exact elimination returns
$S_0,\dots,S_{2\omega-2}$, hence an explicit polynomial left inverse
$S(\lambda)$, verifiable by multiplying the two polynomial matrices. The test is
therefore not merely a yes/no decision.
\end{remark}

\section{Membership in the existential theory of the reals}\label{sec:existsr}

This section places the decision problem at a prescribed order in the existential theory of the reals, and records the time bound that follows from the two algorithms.

\begin{thm}\label{thm:existsr}
For every fixed $q$, $\mathrm{DEC}(q)\in\exists\mathbb R$. Consequently the threshold problem $\mathrm{MINDEC}(Q)$ is in $\exists\mathbb R$ for every $Q$.
\end{thm}

\begin{proof}
Appendix~\ref{app:encoding} constructs, for each $\omega=0,\dots,\min(q,k)$, an existential formula $\Phi_\omega$ with
\begin{equation*}
O(q\omega+\omega k+\omega q+(2\omega-1)\omega q)=O(n^3)
\end{equation*}
real variables and $O(n^3)$ equations of degree at most two, and with no enumeration of column sets. The disjunction
\begin{equation}\label{eq:existsr-disj}
\Phi=\bigvee_{\omega=0}^{\min(q,k)}\Phi_\omega
\end{equation}
has $O(k)$ branches of polynomial size, is built in polynomial time from the input, and its solutions correspond exactly to feasible $L$:
\begin{equation*}
\begin{aligned}
&\Phi\ \text{satisfiable}\\
&\iff\exists\,L=[L_0;R]\ \text{satisfying}\ \eqref{eq:5a}\ \text{and}\ \eqref{eq:5b}.
\end{aligned}
\end{equation*}
Thus $\Phi$ is a polynomial-time reduction to the existential theory of the reals, so $\mathrm{DEC}(q)\in\exists\mathbb R$. Decidability of the target theory is a separate classical result and is not used here. For $\mathrm{MINDEC}(Q)$ take the finite disjunction \eqref{eq:mindec} over $q=r_0,\dots,Q$, which is a polynomial-time reduction of the same kind; hence $\mathrm{MINDEC}(Q)\in\exists\mathbb R$ for every $Q$.
\end{proof}

\begin{cor}\label{cor:time}
For every $r\ge0$, the decision time in the bit model satisfies
\begin{equation*}
T\le\min\bigl\{\mathcal L^{O(q^2)},\
\exp\bigl(O(\ell^2\log(\ell+2))\bigr)\operatorname{poly}(\mathcal L)\bigr\}.
\end{equation*}
\end{cor}

\begin{proof}
The two bounds available are those of Sections~\ref{sec:fixedq} and~\ref{sec:fpt},
\begin{equation*}
\begin{aligned}
T&\le\mathcal L^{O(q^2)}\quad(q\ \text{fixed}),\\
T&\le\exp\bigl(O(\ell^2\log(\ell+2))\bigr)\operatorname{poly}(\mathcal L)
\quad(\ell=\operatorname{rank}(AN)),
\end{aligned}
\end{equation*}
and both algorithms decide the same language, so either may be run and the smaller bound applies.
For $r=0$ the second bound is $\operatorname{poly}(\mathcal L)$ by the test of
Section~\ref{sec:zero}.
\end{proof}


\section{Conclusions}\label{sec:conclusions}

The two upper bounds are parameterized, and the parameter enters a different part of each:
\begin{equation*}
\begin{aligned}
T&\le\mathcal L^{O(q^2)}\quad\text{(exponent governed by }q\text{)},\\
T&\le\exp\bigl(O(\ell^2\log(\ell+2))\bigr)\operatorname{poly}(\mathcal L)\\
&\qquad\text{(factor governed by }\ell=\operatorname{rank}(AN)\text{)}.
\end{aligned}
\end{equation*}
The second parameter does not involve the number $r$ of free rows. The same reduction
settles the neighbouring questions: $\mathrm{OPT}$ is decided by $O(\log\ell)$ calls to the
decision procedure, and $\mathrm{FIND}$ returns exact real algebraic numbers in the arithmetic
model (Corollaries~\ref{cor:opt} and~\ref{cor:find}).
The case of no free rows is decided in $\operatorname{poly}(\mathcal L)$ time by the test of Section~\ref{sec:zero}. Membership in the existential theory of the reals holds for both the prescribed-order and the threshold problem:
\begin{equation*}
\mathrm{DEC}(q)\in\exists\mathbb R,\qquad \mathrm{MINDEC}(Q)\in\exists\mathbb R .
\end{equation*}
The feasibility characterization itself is taken from \cite{fer10,dar25,fer26a,fer26b}; what is new here is the compression to $O(q^2)$ variables, independently of the state dimension, and the complexity bounds that follow from it.

\subsection{Limitations and what is not claimed}\label{sec:limits}

We do not claim that $\mathrm{DEC}\in\mathrm P$, that $\mathrm{DEC}$ is
$\exists\mathbb R$-complete, or any hardness result, and hardness and gadget constructions are
outside the scope of this paper. Neither parameter is bounded in the unrestricted problem, so the
two bounds are compatible with all three outcomes
\begin{equation*}
\begin{aligned}
&\mathrm{DEC}\in\mathrm P;\\
&\mathrm{DEC}\ \text{is}\ \exists\mathbb R\text{-complete};\\
&\mathrm{DEC}\ \text{has an intermediate parameterized classification},
\end{aligned}
\end{equation*}
and the paper makes no claim among them; the unrestricted complexity of $\mathrm{DEC}$ remains
open. The three are stated for $\mathrm{DEC}$ rather than for the minimum-order problem: a
polynomial-time algorithm for $\mathrm{DEC}$ would give one for minimum order, since
Algorithm~\ref{alg:opt} is a binary search over $O(\log n)$ calls to $\mathrm{DEC}$.

Three items remain outside the verified core.
\begin{enumerate}
\item \textbf{The bit complexity of $\mathrm{FIND}$.} Corollary~\ref{cor:find} bounds the
arithmetic operations and the length of the output, but not the bit operations needed to
produce it; the analysis of \cite{bpr06} that we invoke is an arithmetic one, and its Chapter~8
states that bit-length estimates do not follow from it. A bit-model sampling-point theorem would
be needed, and we are not aware of one.
\item \textbf{The $r=0$ branch: resolved, with a residual audit.} For $r=0$ the two conditions reduce to a constant-rank test and to the constancy of the rank of the degree-one pencil $\lambda X-G$. Section~\ref{sec:zero} decides the latter by a finite rational linear system derived from the degree bound $2\omega-2$; the test is deterministic and exact,
\begin{equation*}
\text{unknowns}=\operatorname{poly}(q,\omega),\qquad
\text{equations}=\operatorname{poly}(q,\omega),
\end{equation*}
and it returns an explicit witness. The remaining issue is to complete the explicit bit-complexity audit: the constants in the bit bound of Section~\ref{sec:zero} are standard Hadamard/elimination bounds but have not been checked line by line.
\item \textbf{Coefficient bounds.} The coefficient-size lemma in Section~\ref{sec:fpt} is paper-level; its constants have not been independently audited.
\end{enumerate}

\subsection{The size of the parameter}\label{sec:paramsize}

Two quantities describe the free data: the number of coordinates
\begin{equation}\label{eq:kappa-size}
\kappa_{\rm obs}=r\,d_0\le2rn
\end{equation}
needed to write the $r$ free rows, and the parameter that governs the bound,
\begin{equation}\label{eq:ell-size}
\ell=\operatorname{rank}(AN)\le k\le n .
\end{equation}
The first can be quadratic in $n$; the second cannot. This is the content of the subspace
reduction of Section~\ref{sec:subspace}: a feasible observer is a subspace $W$ of the
compressed row-pair space, its dimension $\omega$ satisfies $\omega+1\le\dim W\le2\omega$ with
$\omega\le\ell$, and the $r$ free rows enter only through the constraint
$\dim(W/W_0)\le r$. The bound of Theorem~\ref{thm:fpt} therefore depends on $\ell$ and not on
$\kappa_{\rm obs}$: the coordinate count \eqref{eq:kappa-size} records how the free rows are
written, not how hard the instance is.

Whether $\ell$ is optimal is open. Two further reductions are conceivable, and they are of
different kinds.

\begin{enumerate}
\item \emph{Quotients of the observer rows.} The conditions depend on row spaces, hence are
invariant under row operations of $L$. Those preserving $L_0$ are block lower-triangular,
\begin{equation*}
T=\begin{bmatrix}I_{r_0}&0\\ U&S\end{bmatrix},\qquad \det S\ne0 ,
\end{equation*}
so two different free blocks can present the same $W$. The reduction of
Section~\ref{sec:subspace} already quotients by this symmetry; the question is whether a
smaller invariant than $\ell$ exists.
\item \emph{Spectral invisibility.} One might hope to divide by directions that change neither
the rank condition nor the finite spectrum. Such a quotient must preserve the interaction
between $H$ and $H_A$, not only $\operatorname{rank}H$; a quotient preserving \eqref{eq:5a} but
changing the spectrum would be invalid. We do not know whether one exists.
\end{enumerate}

\paragraph*{Growth of the parameter-dependent factor}
Renegar's bound is $(sD)^{O(N)}$ with
\begin{equation*}
sD=O(\ell^2),\qquad N=\ell^2+2\ell ,
\end{equation*}
so $f(\ell)=\exp\bigl(O(\ell^2\log(\ell+2))\bigr)$. The result is primarily structural: it
identifies an invariant of the instance, independent of how many free rows are allowed, that
governs the complexity bound, although the doubly exponential dependence on $\ell$ may limit
practical applicability for large values.

\subsection{Reproducibility and evidence}\label{sec:repro}

The results rest on three kinds of evidence, and the accompanying artifact records which statement rests on which.

\begin{enumerate}
\item \emph{Lean~4 with Mathlib.} The statements marked \emph{literal} in
Table~\ref{tab:evidence} are formalized as stated, including the rank identity
\eqref{eq:rankid}, the rank-consistency criterion, the degree-bounded B\'ezout lemma, the
rectangular left inverse with its degree bound, and the sandwich and structural estimates of
Section~\ref{sec:subspace}; those marked \emph{abstract core} are formalized at the level of
linear maps and submodules, with the translation to matrices left to the paper; the single
\emph{ingredient} is one step of the proof of the certificate. Every declaration is checked by the
kernel against only the three standard axioms (\texttt{propext}, \texttt{Classical.choice},
\texttt{Quot.sound}). The development is not a formal verification of the algorithmic theorems
themselves: the reductions, the certificate and the complexity accounting are paper-level, and the
formal statements cover the algebraic and linear-algebraic components listed in the table.
\item \emph{Paper-level proofs.} The Cauchy--Binet compression of the certificate, the coefficient-size estimates and the existential-real encoding.
\item \emph{Exact rational arithmetic.} The identities on which the constructions rest are re-checked, both on randomly generated instances and on the constructed families; this excludes arithmetic slips but does not replace the proofs.
\end{enumerate}

The Lean development and the exact-arithmetic scripts form the accompanying artifact and are kept at \url{https://github.com/Zhengshuji/q1-functional-observer}. Table~\ref{tab:evidence} gives the correspondence. The matrix-level identifications that the Lean core does not cover, namely the pencil/gcd bridge, the matrix form of $\mathcal T$, the converse step from a polynomial left inverse to the linear system \eqref{eq:R0}, and the encoding, are paper-level and are the ones re-checked by exact arithmetic. Of the classical facts invoked by citation above, the rank identity of Lemma~\ref{lem:rankid}, the solvability criterion of Lemma~\ref{lem:rankcons} and the coprimality and degree-bounded B\'ezout steps of Lemma~\ref{lem:bezout} are also proved in the artifact, so for those three the citation serves only to keep the paper's proofs short; the remaining cited facts are used as references and are not formalized.

\begin{table*}[!t]
\centering
\caption{Where each statement of the paper is verified. \emph{Literal}: the Lean statement is the
statement of the paper. \emph{Abstract core}: the Lean statement is the linear-map or submodule
version, and the translation to matrices is paper-level. \emph{Ingredient}: one step of the proof,
not the statement. Statements not listed are paper-level throughout.}
\label{tab:evidence}
\footnotesize
\begin{tabular}{@{}lll@{}}
\toprule
Statement & Lean evidence & Formal name \\
\midrule
Rank identity (Lemma~\ref{lem:rankid}), \eqref{eq:rankid} & literal &
\texttt{rank\_vstack\_add\_rank\_mul} \\
Compression criterion (Lemma~\ref{lem:compression}) & abstract core &
\texttt{ker\_compressed\_iff} \\
Space of row pairs (Lemma~\ref{lem:attainable}) & abstract core & \texttt{freePairSpace},
\texttt{kappaObs\_eq\_mul\_freePairSpace} \\
Sandwich and realizability (Lemma~\ref{lem:sandwich}) & literal &
\texttt{rowSpan\_fromRows\_left}, \texttt{mem\_rowSpan\_iff\_exists\_vecMul} \\
Conditions depend on $W$ (Lemma~\ref{lem:W-conditions}) & abstract core &
\texttt{rank\_fromRows\_eq\_rank\_iff}, \texttt{rowSpace\_eq} \\
Structural bound, upper half (Lemma~\ref{lem:struct}) & literal &
\texttt{finrank\_le\_two\_mul\_of\_pi1\_eq\_pi2} \\
Rank consistency (Lemma~\ref{lem:rankcons}) & literal &
\texttt{exists\_mulVec\_eq\_iff\_rank\_gram} \\
B\'ezout with degree bound (Lemma~\ref{lem:bezout}) & literal & \texttt{exists\_bezout\_family} \\
Rectangular left inverse (Proposition~\ref{prop:zero}) & literal &
\texttt{exists\_leftInverse\_linear} \\
Gram certificate, kernel identity (Proposition~\ref{prop:gram}) & ingredient &
\texttt{sum\_sq\_zero\_iff} \\
Coefficient system \eqref{eq:R0} $\Rightarrow$ left inverse & literal (one direction) &
\texttt{spoly\_mul\_pencil\_eq\_one} \\
Matrix-product coefficients & literal & \texttt{coeff\_mul\_apply} \\
Degree bounds $\deg\det$, $\deg\operatorname{adj}$ & literal &
\texttt{natDegree\_det\_le}, \texttt{natDegree\_adjugate\_le} \\
Zero-free-row test (Section~\ref{sec:zero}) & paper-level & --- \\
Encoding (Theorem~\ref{thm:existsr}) & paper-level & --- \\
\bottomrule
\end{tabular}
\end{table*}

\appendices
\section{The existential-real encoding}\label{app:encoding}

This appendix constructs the encoding used in the proof of Theorem~\ref{thm:existsr}. Throughout, $q$ is fixed and the notation is that of Section~\ref{sec:compressed}: $N$ is a basis matrix of $\ker[C;CA]$, $H=LN$, $H_A=LAN$, and $\omega=\operatorname{rank}H$.

For fixed $q$, \eqref{eq:5b} is the condition that the $q\times\omega$ pencil
\begin{equation}
Q(\lambda)=\lambda X-G
\end{equation}
has rank $\omega$ for every $\lambda\in\mathbb C$. This is stronger than having rank $\omega$ over the fraction field $\mathbb C(\lambda)$: the latter does not give $\operatorname{rank}Q(\lambda)=\omega$ at every $\lambda$, as
\begin{equation*}
Q(\lambda)=\begin{bmatrix}\lambda\\ \lambda\end{bmatrix},\qquad
\operatorname{rank}_{\mathbb C(\lambda)}Q=1,\qquad \operatorname{rank}Q(0)=0
\end{equation*}
shows.

By the Quillen--Suslin/B\'ezout criterion for polynomial matrices over a principal ideal
domain \cite{fab90,lan02}, $Q$ has rank $\omega$ at every $\lambda$ iff its maximal minors generate the unit ideal in $\mathbb C[\lambda]$, equivalently iff there is a polynomial left inverse,
\begin{equation*}
\begin{aligned}
&\operatorname{rank}Q(\lambda)=\omega\ \text{ for all }\lambda\\
&\iff\langle f_I\rangle=\mathbb C[\lambda]\\
&\iff\exists\,S\in\mathbb C[\lambda]^{\omega\times q}\ \text{with}\ S(\lambda)Q(\lambda)=I_\omega .
\end{aligned}
\end{equation*}

The degree bound is $2\omega-2$. For each $\omega$-element row set $I\subseteq\{1,\dots,q\}$, let $Q_I=E_IQ$ be the corresponding $\omega\times\omega$ submatrix and $f_I=\det Q_I$. If $Q$ has rank $\omega$ at every $\lambda$, the $f_I$ have gcd one, and Lemma~\ref{lem:bezout} gives
\begin{equation}
\sum_I g_I f_I=1,\qquad \deg g_I<\omega .
\end{equation}
Let $S_I$ be the $\omega\times q$ matrix whose $I$-columns are $\operatorname{adj}(Q_I)$ and whose other columns are zero. Then
\begin{equation*}
S_IQ=\operatorname{adj}(Q_I)E_IQ=f_II_\omega,
\end{equation*}
so
\begin{equation*}
S=\sum_I g_I S_I\ \Longrightarrow\ SQ=I_\omega,\qquad
\deg S\le2\omega-2,
\end{equation*}
because $\deg g_I<\omega$ and $\deg\operatorname{adj}(Q_I)\le\omega-1$.

For the $\exists\mathbb R$ encoding the minors are not enumerated. For each $\omega=0,\dots,\min(q,k)$ introduce the blocks
\begin{equation*}
\begin{gathered}
X\in\mathbb R^{q\times\omega},\qquad B\in\mathbb R^{\omega\times k},\\
G\in\mathbb R^{q\times\omega},\qquad U\in\mathbb R^{\omega\times q},\\
V\in\mathbb R^{k\times\omega},
\end{gathered}
\end{equation*}
and the coefficients of $S(\lambda)=\sum_{j=0}^{2\omega-2}S_j\lambda^j$, where $S_j\in\mathbb R^{\omega\times q}$, and impose
\begin{equation*}
\begin{gathered}
H=XB,\qquad H_A=GB,\\
UX=I_\omega,\qquad BV=I_\omega,\\
S(\lambda)(\lambda X-G)=I_\omega .
\end{gathered}
\end{equation*}
The polynomial identity is linear in the coefficients $S_j$; comparing coefficients of $\lambda^j$ makes it the linear system \eqref{eq:R0}. The rank conditions $UX=I_\omega$ and $BV=I_\omega$ are polynomial equations, and $BV=I_\omega$ makes $B$ full row rank, hence
\begin{equation}
\operatorname{rank}P(\lambda)=\operatorname{rank}((\lambda X-G)B)
=\operatorname{rank}(\lambda X-G)=\omega
\end{equation}
for all $\lambda$. Also $H_A=GB=(GU)XB=FH$ with $F=GU$, so \eqref{eq:5a} holds.

It remains to tie $(H,H_A)$ to the input. The rows of $H=LN$, $H_A=LAN$ satisfy, with $h_i^0=(L_0N)_i$ and $g_i^0=(L_0AN)_i$,
\begin{equation}
X_iB=h_i^0,\qquad G_iB=g_i^0\qquad(1\le i\le r_0),
\tag{B.1}\label{eq:B1}
\end{equation}
and, for the free rows, $(H_i,H_{A,i})\in\mathcal T=\operatorname{row}[N\ AN]$, equivalently
\begin{equation}
D\begin{bmatrix}(X_iB)^{\mathsf T}\\ (G_iB)^{\mathsf T}\end{bmatrix}=0
\qquad(r_0<i\le q),
\tag{B.2}\label{eq:B2}
\end{equation}
where $D$ is a rational matrix with $\ker D=\mathcal T$. Both are polynomial (in fact bilinear) equations. They are not optional: without \eqref{eq:B1}--\eqref{eq:B2} the encoding would accept arbitrary pairs $(H,H_A)$.

The correctness splits into two independent statements.

\paragraph*{Soundness}
Suppose the system has a solution. Set $L$ with first $r_0$ rows $L_0$ and, for each $i>r_0$, a preimage $r_i$ with $(X_iB,G_iB)=r_i[N\ AN]$, which exists by \eqref{eq:B2}. Then
\begin{equation*}
\begin{aligned}
H&=LN,\qquad H_A=LAN,\\
H_A&=GB=(GU)XB=FH,\qquad F=GU,
\end{aligned}
\end{equation*}
so \eqref{eq:5a} holds; and
\begin{equation*}
\operatorname{rank}(\lambda H-H_A)=\operatorname{rank}(\lambda X-G)=\omega
\end{equation*}
for all $\lambda$, which gives \eqref{eq:5b} and uses $BV=I_\omega$ (full row rank of $B$) together with $S(\lambda)(\lambda X-G)=I_\omega$.

\paragraph*{Completeness}
A feasible $L$ supplies $X,B,U,V$ from a rank factorization $H=XB$ with $B_{:,J}=I_\omega$, $G=FX$ with $H_A=FH$ (Lemma~\ref{lem:compression}), and $S$ from the left-inverse construction above; then \eqref{eq:B1} and \eqref{eq:B2} hold by construction, and the remaining equations are identities.

The two constraints $UX=I_\omega$ and $BV=I_\omega$ are not redundant bookkeeping: the first produces the factorization $H_A=FH$, the second makes $\operatorname{rank}((\lambda X-G)B)=\operatorname{rank}(\lambda X-G)$ hold at every $\lambda$.

The $\omega=0$ case is the branch $H=H_A=0$, which \eqref{eq:B1} forces when $L_0N=0$ and rejects otherwise; it is a separate algebraic branch, not a degree-negative instance of \eqref{eq:R0}. The number of variables is
\begin{equation}
O(q\omega+\omega k+q\omega+\omega q+k\omega+(2\omega-1)\omega q)
=O(n^3),
\end{equation}
the number of equations is polynomial, and the degree of the left inverse is $2\omega-2\le2n$. The disjunction \eqref{eq:existsr-disj} over $\omega$ has at most $n+1$ branches. Hence the encoding has polynomial size. This proves $\mathrm{DEC}(q)\in\exists\mathbb R$ for every fixed $q$, and hence $\mathrm{MINDEC}(Q)\in\exists\mathbb R$.

\section{A worked example}\label{app:example}

We run the procedure on the instance of \cite[Sec.~VII]{fer26b}, which has $n=8$, $p=2$ and
\begin{equation*}
\begin{gathered}
C=\begin{bmatrix}-2&-2&2&-2&-1&0&1&-3\\ 0&0&1&2&2&3&-1&-2\end{bmatrix},\\
L_0=\begin{bmatrix}-2&-1&-2&-3&-1&3&2&0\\ 1&-1&-2&1&0&2&0&1\end{bmatrix},
\end{gathered}
\end{equation*}
and
\begin{equation*}
A=\begin{bmatrix}
0&1&0&1&0&0&0&0\\
0&-1&1&0&0&0&0&-1\\
1&0&-1&1&0&0&1&-1\\
-1&-1&0&-2&0&0&0&0\\
0&0&0&0&-2&0&0&1\\
0&0&0&0&0&-2&1&0\\
0&0&0&0&0&-1&-2&1\\
0&0&0&0&0&0&1&-2
\end{bmatrix}.
\end{equation*}
The compression of Section~\ref{sec:dimension} gives
\begin{equation*}
\begin{gathered}
k=\dim\ker[C;CA]=4,\qquad d_0=\operatorname{rank}[N\ AN]=6,\\
s_0=\dim\operatorname{row}(L_0M)=2,
\end{gathered}
\end{equation*}
and, in the notation of \eqref{eq:ell},
\begin{equation*}
S=\operatorname{row}(AN)=\mathbb R^4,\qquad \ell=4,\qquad m'=4 .
\end{equation*}
Here $\ell=k$, so the subspace reduction is vacuous and the content of the example is the
certificate alone. The coordinate count is $\kappa_{\rm obs}=rd_0=12$ already at $r=2$, whereas
the parameter of Theorem~\ref{thm:fpt} is $\ell=4$.

At $t=2$, that is at order $q=r_0+t=4$, the procedure of Section~\ref{sec:fpt} accepts. An explicit
witness is the four-row observer matrix $L=[L_0;R]$ with
\begin{equation*}
R=\begin{bmatrix}
-\tfrac{42919}{32067} & -\tfrac{55096}{10689} & 1 & 0 & 0 & 0 & 0 & 0\\
\tfrac{3283}{4581} & \tfrac{118}{1527} & 0 & 1 & 0 & 0 & 0 & 0
\end{bmatrix}.
\end{equation*}
For this matrix,
\begin{equation*}
\operatorname{rank}[L;CA;C]=7=\operatorname{rank}[LA;L;CA;C],
\end{equation*}
which is \eqref{eq:5a}. The pencil $[\lambda L-LA;\,CA;\,C]$ is $8\times8$; all $64$ of its maximal
minors are nonzero and their gcd is $1$, while every larger minor vanishes, so its rank is $7$ for
every $\lambda\in\mathbb C$, which is \eqref{eq:5b}. Hence $\mathrm{DEC}(4)$ holds and
$q_{\min}\le4$ for this instance.

We record that \cite[Sec.~VII-B]{fer26b} reports this instance as infeasible at order $4$, and
hence $\{5,6\}$ as its order spectrum. The matrix $L$ above is a four-row witness satisfying the
two conditions, so \emph{that report is in error}. The error is not in the characterization but in
the way it is evaluated there: the condition-(A) family is parameterized chart by chart and a
branch of it is missed. The formulation of Section~\ref{sec:fpt} never forms that family---it
quantifies over the chart coordinates and the certificate jointly---which is why it cannot miss a
branch. Nothing else in this paper depends on the discrepancy: the characterization of
\cite{fer26b} is used only for the existence of a nullspace matrix, and the witness supplies one.

\paragraph*{The two procedures on this instance}
Write $M=[N\ AN]$ and, for the accepted $L$, $H=LN$, $H_A=LAN$.

\emph{OPT.} Algorithm~\ref{alg:opt} with $Q=6$ evaluates $\mathrm{DEC}(6)$, which is feasible, so the
answer is finite; the search interval is $[0,\min(4,2\ell)]=[0,4]$. The call at $t=2$, that is at
order $4$, is the certificate above and accepts. The call at $t=0$, that is at $L=L_0$, rejects: for
$L_0$ we have $\operatorname{rank}[L_0;CA;C]=6$ but $\operatorname{rank}[L_0A;L_0;CA;C]=7$, so
\eqref{eq:5a} fails. At $t=1$, that is at order $3$, no feasible $W$ exists, and Proposition~\ref{prop:lowerbound}
decides it without any spectral test. The two blocks of $W_0$ are $U_0=L_0N$ and $V_0=L_0AN$, with
\begin{equation*}
\operatorname{rank}U_0=2,\qquad \operatorname{rank}[\,U_0;\,V_0\,]=3,
\end{equation*}
so \eqref{eq:qmin-lb} gives $t\ge3-2+1=2$: every feasible $W$ has quotient dimension at least $2$.
Concretely, a candidate with $t=1$ would have $\dim W=s_0+1=3$, so Lemma~\ref{lem:struct} would force
$\omega=2$, and $\pi_1(W)=\pi_2(W)=U$ with $\dim U=\dim U_0=2$ would force $U=U_0$ and hence
$V_0\subseteq U_0$, which fails. Therefore $\mathrm{DEC}(3)$ holds for no candidate, $\mathrm{DEC}(4)$
holds by the certificate above, and

\begin{equation*}
q_{\min}=4
\end{equation*}

for this instance. The value $5$ reported in \cite[Sec.~VII]{fer26b} is therefore not $q_{\min}$.

\emph{FIND.} For the accepted $L$ the subspace $W=\operatorname{row}[H\ H_A]$ is four-dimensional and
$W_0=\operatorname{row}[H_0\ H_{A,0}]$ is two-dimensional, so $t=\dim(W/W_0)=2$; moreover $W_0\subseteq
W$, and the rows $r_1M$ and $r_2M$ of the two free rows, together with the two rows of $W_0$, span
$W$:
\begin{equation*}
\begin{gathered}
\dim W=4,\qquad \dim W_0=2,\\
\operatorname{rank}[\,W_0;\,r_1M;\,r_2M\,]=4 .
\end{gathered}
\end{equation*}
Step~2 of Procedure~\ref{proc:find} is therefore realized by $u_i=r_i$, and step~3 pads with
$Q-r_0-t=6-2-2=2$ zero rows. Only step~1 of the procedure is not exercised here, because a witness
$W$ was found directly rather than by solving the algebraic system.

The two exact rank computations of $U_0$ and $[U_0;V_0]$ are the whole of the argument at order
$3$; they are reproduced by the script archived with the artifact, together with the rank and
minor computations used above.

\begin{thebibliography}{99}
\bibitem{fer10} J. A. Fernando, H. Trinh, L. Jennings. Functional observability and the design of minimum order linear functional observers. \emph{IEEE Trans. Automat. Control} 55(5):1268--1273, 2010.
\bibitem{dar25} M. Darouach, J. A. Fernando. On functional observability and functional observer design. \emph{Automatica}, vol. 173, Art. no. 112115, 2025, doi: 10.1016/j.automatica.2025.112115.
\bibitem{fer26a} J. A. Fernando. A minimum-order functional observer beyond the Darouach and Luenberger constructions. arXiv:2609.33615, 2026 (preprint).
\bibitem{fer26b} J. A. Fernando. Complete characterization of minimum-order functional observers from Darouach to Luenberger. arXiv:2609.37992, 2026 (preprint).
\bibitem{zha25} Y. Zhang, T. Fernando, M. Darouach. Functional
  observability, structural functional observability, and optimal sensor
  placement. \emph{IEEE Trans. Automat. Control}, vol. 70, no. 3,
  pp. 1592--1607, 2025, doi: 10.1109/TAC.2024.3462558.
  (arXiv:2307.08923.)
\bibitem{dey21} P. Dey, N. Balachandran, D. Chatterjee. Complexity of
  constrained sensor placement problems for optimal observability.
  \emph{Automatica}, vol. 131, Art. no. 109758, 2021,
  doi: 10.1016/j.automatica.2021.109758.
\bibitem{lof26} J. L\"{o}fberg. Static output-feedback stabilization is
  NP-hard. arXiv:2609.16886, 2026 (preprint).
\bibitem{bar26} G. Barkai, I. Shames. On the NP-hardness of unconstrained
  static output feedback stabilization. arXiv:2609.38728, 2026 (preprint).
\bibitem{liu26} Z. Liu. Computational complexity of static output feedback
  stabilization. arXiv:2610.00390, 2026 (preprint).
\bibitem{maj26} A. Majumdar. Affine rank minimization is
  $\exists\mathbb{R}$ complete. arXiv:2602.14037, 2026 (preprint).
\bibitem{hj13} R. A. Horn, C. R. Johnson. \emph{Matrix Analysis}. 2nd ed.
Cambridge University Press, 2013.
\bibitem{vzgg13} J. von zur Gathen, J. Gerhard. \emph{Modern Computer Algebra}.
3rd ed. Cambridge University Press, 2013.
\bibitem{emi99} I. Z. Emiris, B. Mourrain. Matrices in elimination theory. \emph{J. Symbolic Comput.} 28(1--2):3--44, 1999, doi: 10.1006/jsco.1998.0266.
\bibitem{bpr06} S. Basu, R. Pollack, M.-F. Roy. \emph{Algorithms in Real
  Algebraic Geometry}, 2nd ed. Berlin, Germany: Springer, 2006,
  Algorithms and Computation in Mathematics, vol. 10.
\bibitem{ren92} J. Renegar. On the computational complexity and geometry of the first-order theory of the reals. Part II: The general decision problem. \emph{J. Symbolic Comput.} 13(3):301--327, 1992, doi: 10.1016/S0747-7171(10)80004-5.
\bibitem{fab90} A. Fabiano, D. C. Struppa. Bounds for the degrees of the entries of left
inverses of polynomial matrices. \emph{Multidimensional Systems and Signal Processing},
vol.~1, no.~1, pp.~67--74, 1990, doi: 10.1007/BF01812207.
\bibitem{bee88} Th. Beelen, P. Van Dooren. An improved algorithm for the computation of Kronecker's canonical form of a singular pencil. \emph{Linear Algebra Appl.} 105:9--65, 1988.
\bibitem{lan02} S. Lang. \emph{Algebra}. Revised 3rd ed.
Graduate Texts in Mathematics 211, Springer, 2002.
\end{thebibliography}

\end{document}
